\chaptermark{Publicación excepcional del continuo}
\label{chap:83-20}

\input{ampliaciones_sucesoras_20260824/parte_i_ii/tex/bloques/salto_espectral_4_2_6_54}

% Unidad U018. Paley, extension cuadratica finita y codigo ternario.
% Redaccion autonoma desde la autoridad material 1063.

\section{Forma de conferencia de Paley, extensión cuadrática finita y código ternario autodual}
\label{p12-83-c20-s1:chap:paley-extension-cuadratica-codigo-ternario}

La tétrada \(B_K\) construida en la unidad anterior todavía no posee una
estrella de Witt: para definirla es preciso construir primero el sistema de
incidencia que contiene sus hexadas. Esta unidad realiza ese paso desde el
calibre proyectivo de las seis unidades módulo nueve, obtiene una raíz
cuadrática interna de \(-I_6\), y construye el código ternario gráfico de
longitud doce. Ninguna identificación se deduce de una mera igualdad de
cardinales.

\subsection{Calibre proyectivo de las \(6\) unidades módulo \(9\)}

Sea
\[
 \mathcal U_9=(\mathbb Z/9\mathbb Z)^\times
 =\{1,2,4,8,7,5\},
\]
ordenado por las potencias sucesivas de \(2\) módulo \(9\). Fijamos la
biyección de coordenadas
\begin{equation}
 \theta:\mathcal U_9\longrightarrow\mathbb P^1(\mathbb F_5),
 \qquad
 \begin{array}{c|cccccc}
 u&1&2&4&8&7&5\\ \hline
 \theta(u)&\infty&0&1&2&3&4
 \end{array}.
 \label{p12-83-c20-s1:eq:calibre-unidades-paley}
\end{equation}
La aplicación \(\theta\) es un calibre de coordenadas. No se afirma que
sea un homomorfismo entre la multiplicación de \(\mathcal U_9\) y una
operación sobre \(\mathbb P^1(\mathbb F_5)\).

Sea
\[
 \chi:\mathbb F_5\longrightarrow\{-1,0,1\}
\]
el carácter cuadrático:
\[
 \chi(0)=0,
 \qquad
 \chi(1)=\chi(4)=1,
 \qquad
 \chi(2)=\chi(3)=-1.
\]
En el orden \((\infty,0,1,2,3,4)\), definimos
\begin{align}
 (C_P)_{\infty,x}&=(C_P)_{x,\infty}=1,
 &&x\in\mathbb F_5,\label{p12-83-c20-s1:eq:paley-fila-infinito}\\
 (C_P)_{x,y}&=\chi(x-y),
 &&x,y\in\mathbb F_5, x\neq y,\label{p12-83-c20-s1:eq:paley-entradas-finitas}\\
 (C_P)_{x,x}&=0.
 \label{p12-83-c20-s1:eq:paley-diagonal}
\end{align}
Así,
\begin{equation}
 C_P=
 \begin{pmatrix}
 0& 1& 1& 1& 1& 1\\
 1& 0& 1&-1&-1& 1\\
 1& 1& 0& 1&-1&-1\\
 1&-1& 1& 0& 1&-1\\
 1&-1&-1& 1& 0& 1\\
 1& 1&-1&-1& 1& 0
 \end{pmatrix}.
 \label{p12-83-c20-s1:eq:matriz-conferencia-paley-autonoma}
\end{equation}

\begin{lemma}[Suma de correlación cuadrática]
\label{p12-83-c20-s1:lem:correlacion-cuadratica-f5}
Para \(a\neq b\) en \(\mathbb F_5\),
\begin{equation}
 \sum_{t\in\mathbb F_5}
 \chi(t-a)\chi(t-b)=-1.
 \label{p12-83-c20-s1:eq:correlacion-cuadratica-f5}
\end{equation}
\end{lemma}

\begin{proof}
La sustitución \(u=(t-a)/(b-a)\) es una biyección de \(\mathbb F_5\).
Como \(\chi((b-a)^2)=1\), la suma se reduce a
\[
 \sum_{u\in\mathbb F_5}\chi(u(u-1)).
\]
La evaluación de \(u=0,1,2,3,4\) da respectivamente
\(0,0,1,-1,-1\), cuya suma es \(-1\).
\end{proof}

\begin{proposition}[Identidad de conferencia]
\label{p12-83-c20-s1:prop:paley-identidad-conferencia-autonoma}
La matriz \(C_P\) es simétrica y satisface
\begin{equation}
 \boxed{C_PC_P^{\mathsf T}=5I_6.}
 \label{p12-83-c20-s1:eq:paley-identidad-conferencia-autonoma}
\end{equation}
\end{proposition}

\begin{proof}
Cada fila contiene cinco entradas de módulo uno y una entrada nula; su
norma cuadrada es cinco. Para dos filas finitas distintas, la contribución
de la coordenada \(\infty\) es \(+1\), mientras
\cref{p12-83-c20-s1:lem:correlacion-cuadratica-f5} da \(-1\) para las cinco coordenadas
finitas. Su producto escalar es cero. La suma de un carácter no trivial
sobre \(\mathbb F_5\) es cero, por lo que la fila correspondiente a
\(\infty\) también es ortogonal a cada fila finita.
\end{proof}

\subsection{Reducción ternaria y raíz cuadrática interna de \texorpdfstring{\(-I_6\)}{-I₆}}

Reducimos \(C_P\) módulo tres, con \(-1\equiv2\pmod3\):
\begin{equation}
 A_W=
 \begin{pmatrix}
 0&1&1&1&1&1\\
 1&0&1&2&2&1\\
 1&1&0&1&2&2\\
 1&2&1&0&1&2\\
 1&2&2&1&0&1\\
 1&1&2&2&1&0
 \end{pmatrix}
 \in M_6(\mathbb F_3).
 \label{p12-83-c20-s1:eq:matriz-paley-witt-ternaria}
\end{equation}

\begin{theorem}[Raíz cuadrática interna de \(-I_6\)]
\label{p12-83-c20-s1:thm:aw-cuadrado-menos-identidad-autonomo}
Se cumple
\begin{equation}
 \boxed{A_W^2=-I_6.}
 \label{p12-83-c20-s1:eq:aw-cuadrado-menos-identidad-autonomo}
\end{equation}
En particular,
\[
 A_W^{-1}=-A_W.
\]
\end{theorem}

\begin{proof}
La \cref{p12-83-c20-s1:prop:paley-identidad-conferencia-autonoma} da
\[
 A_WA_W^{\mathsf T}
 \equiv5I_6
 \equiv2I_6
 =-I_6
 \pmod3.
\]
Como \(A_W=A_W^{\mathsf T}\), se obtiene \(A_W^2=-I_6\).
\end{proof}

La identidad anterior es la residencia algebraica de la raíz de \(-1\)
en este corredor. No procede del factor \(1/4\) de la inversión Hadamard:
aquél es un escalar racional forzado por \(H_4^2=4I_4\); ésta es una
estructura cuadrática sobre característica tres.

\subsection{Extensión cuadrática a \texorpdfstring{\(\mathbb F_9\)}{F₉}}

Definimos
\begin{equation}
 \mathbb F_9=\mathbb F_3[\iota]/(\iota^2+1),
 \qquad \iota^2=-1.
 \label{p12-83-c20-s1:eq:definicion-f9-autonoma}
\end{equation}
El polinomio \(X^2+1\) no tiene raíces en \(\mathbb F_3\), por lo que el
cociente es un cuerpo de nueve elementos. Para
\(V=\mathbb F_3^6\), sea
\[
 V_9=V\otimes_{\mathbb F_3}\mathbb F_9.
\]
Los operadores
\begin{equation}
 P_\iota=\frac12(I-\iota A_W),
 \qquad
 P_{-\iota}=\frac12(I+\iota A_W)
 \label{p12-83-c20-s1:eq:proyectores-f9-paley}
\end{equation}
están bien definidos porque \(2\) es invertible en \(\mathbb F_3\).

\begin{proposition}[Proyectores conjugados]
\label{p12-83-c20-s1:prop:proyectores-conjugados-paley}
Se cumplen
\[
 P_\iota^2=P_\iota,
 \qquad
 P_{-\iota}^2=P_{-\iota},
 \qquad
 P_\iota P_{-\iota}=0,
 \qquad
 P_\iota+P_{-\iota}=I.
\]
\end{proposition}

\begin{proof}
Usando \(\iota^2=-1\) y \(A_W^2=-I\),
\[
 (\iota A_W)^2=I.
\]
Las cuatro identidades son entonces las identidades elementales de los
proyectores \((I\mp\iota A_W)/2\) de una involución.
\end{proof}

\begin{theorem}[Descomposición espectral conjugada]
\label{p12-83-c20-s1:thm:descomposicion-espectral-f9-paley}
Sea
\[
 E_{\pm\iota}=\ker(A_W\mp\iota I).
\]
Entonces
\begin{equation}
 V_9=E_\iota\oplus E_{-\iota},
 \qquad
 \dim_{\mathbb F_9}E_\iota
 =\dim_{\mathbb F_9}E_{-\iota}=3.
 \label{p12-83-c20-s1:eq:descomposicion-espectral-f9-paley}
\end{equation}
\end{theorem}

\begin{proof}
El polinomio mínimo de \(A_W\) divide \(X^2+1\), que se escinde sobre
\(\mathbb F_9\) con raíces simples \(\pm\iota\). Los proyectores de
\cref{p12-83-c20-s1:prop:proyectores-conjugados-paley} producen la suma directa. Como
\(A_W\) tiene coeficientes en \(\mathbb F_3\), el automorfismo de
Frobenius intercambia los dos autoespacios; sus dimensiones son iguales y
suman seis.
\end{proof}

Esta construcción dota \(\mathbb F_3^6\) de estructura de módulo de rango
tres sobre \(\mathbb F_9\), mediante
\begin{equation}
 (a+b\iota)\cdot v:=av+b(vA_W).
 \label{p12-83-c20-s1:eq:accion-f9-sobre-f3-seis}
\end{equation}
La igualdad \(3^6=9^3\) queda así acompañada por una acción explícita; no
se usa como argumento cardinal.

\subsection{Código gráfico ternario de longitud \(12\)}

Definimos
\begin{equation}
 c:\mathbb F_3^6\longrightarrow\mathbb F_3^{12},
 \qquad
 c(q)=(q,qA_W),
 \label{p12-83-c20-s1:eq:aplicacion-codigo-grafico-witt}
\end{equation}
y
\begin{equation}
 C_W:=\operatorname{im}c.
 \label{p12-83-c20-s1:eq:definicion-codigo-witt-ternario}
\end{equation}
Su matriz generadora es
\[
 G_W=(I_6\mid A_W).
\]

\begin{theorem}[Inyectividad y autodualidad]
\label{p12-83-c20-s1:thm:cw-inyectivo-autodual}
La aplicación \(c\) es inyectiva, \(\dim C_W=6\) y
\begin{equation}
 \boxed{C_W=C_W^\perp.}
 \label{p12-83-c20-s1:eq:cw-autodual}
\end{equation}
\end{theorem}

\begin{proof}
La primera mitad de \(c(q)\) es \(q\), por lo que
\(\operatorname{pr}_1\circ c=\operatorname{id}\); esto prueba la
inyectividad y la dimensión. Además,
\[
 G_WG_W^{\mathsf T}
 =I_6+A_WA_W^{\mathsf T}
 =I_6+2I_6=0
\]
en \(\mathbb F_3\). Por tanto \(C_W\subseteq C_W^\perp\). Ambos espacios
tienen dimensión seis, luego son iguales.
\end{proof}

Para \(r=0,\ldots,12\), sea
\[
 N_r:=\#\{q\in\mathbb F_3^6:
            \operatorname{wt}(q,qA_W)=r\}.
\]
La enumeración de las \(3^6=729\) entradas produce
\begin{equation}
 N_0=1,
 \qquad
 N_6=264,
 \qquad
 N_9=440,
 \qquad
 N_{12}=24,
 \label{p12-83-c20-s1:eq:censo-pesos-cw}
\end{equation}
y \(N_r=0\) en los demás pesos.

\begin{theorem}[Enumerador y distancia mínima]
\label{p12-83-c20-s1:thm:enumerador-cw-autonomo}
El código \(C_W\) posee parámetros
\[
 \boxed{C_W=[12,6,6]_3}
\]
y enumerador de pesos
\begin{equation}
 \boxed{W_{C_W}(y)=1+264y^6+440y^9+24y^{12}.}
 \label{p12-83-c20-s1:eq:enumerador-pesos-cw}
\end{equation}
\end{theorem}

\begin{proof}
La dimensión y la autodualidad ya se demostraron. El censo
\eqref{p12-83-c20-s1:eq:censo-pesos-cw} recorre las \(729\) palabras sin muestreo. Su
menor peso no nulo es seis, que es la distancia mínima de un código lineal.
La suma \(1+264+440+24=729\) controla que ninguna palabra fue omitida.
\end{proof}

\subsection{Naturaleza de la estructura obtenida}

Se han construido tres objetos distintos y tres flechas tipadas:
\[
 \mathcal U_9
 \xrightarrow{\theta}\mathbb P^1(\mathbb F_5)
 \xrightarrow{\chi}C_P
 \xrightarrow{\bmod3}A_W,
\]
\[
 A_W^2=-I_6
 \quad\Longrightarrow\quad
 \mathbb F_3^6\cong\mathbb F_9^3
 \quad\Longrightarrow\quad
 C_W=\{(q,qA_W):q\in\mathbb F_3^6\}.
\]
La primera flecha es un calibre, la segunda usa el carácter cuadrático y la
tercera es reducción de coeficientes. La estructura sobre
\(\mathbb F_9\) es finita; no se identifica con una estructura compleja
real o analítica.

\subsection{Obstrucciones del corredor Paley–\(C_W\): biyección, \(A_W^2=-I_6\), autodualidad y censo de \(729\) palabras}

\begin{criterion}[Refutación del corredor Paley--\(C_W\)]
La construcción queda refutada si ocurre cualquiera de los hechos
siguientes:
\begin{enumerate}
 \item \(\theta\) no es biyectiva o se usa como homomorfismo sin prueba;
 \item la suma \eqref{p12-83-c20-s1:eq:correlacion-cuadratica-f5} difiere de \(-1\);
 \item \(C_PC_P^{\mathsf T}\neq5I_6\);
 \item \(A_W^2\neq-I_6\);
 \item los proyectores \eqref{p12-83-c20-s1:eq:proyectores-f9-paley} no son
       complementarios;
 \item la aplicación \(q\mapsto(q,qA_W)\) pierde inyectividad o
       autodualidad;
 \item el censo no suma \(729\), difiere de
       \eqref{p12-83-c20-s1:eq:censo-pesos-cw}, o presenta una palabra de peso entre uno
       y cinco;
 \item se infiere \(\mathbb F_3^6\cong\mathbb F_9^3\) únicamente de
       \(3^6=9^3\), omitiendo la acción \eqref{p12-83-c20-s1:eq:accion-f9-sobre-f3-seis}.
\end{enumerate}
\end{criterion}


% Unidad U019. Sistema de Witt, campo de estrellas y grupo M12.
% Redaccion autonoma desde la autoridad material 1063.

\section{Sistema de Steiner ternario, involuciones de estrella de Witt y generación de \texorpdfstring{\(M_{12}\)}{M12}}
\label{p12-83-c20-s2:chap:sistema-witt-estrellas-m12}

La unidad precedente construyó el código autodual
\(C_W=[12,6,6]_3\). Ahora se proyectivizan sus palabras de peso seis,
se prueba la propiedad de Steiner, y sólo entonces se define la estrella
asociada a la tétrada \(B_K\). Este orden impide convertir una figura de
incidencia en un objeto añadido por semejanza visual.

\subsection{Soportes proyectivos de peso \(6\)}

Sea
\begin{equation}
 \mathcal H_W
 :=\{\operatorname{supp}(c):c\in C_W,
                              \operatorname{wt}(c)=6\}.
 \label{p12-83-c20-s2:eq:familia-soportes-peso-seis}
\end{equation}
Cada soporte aparece para el par de palabras \(\{c,-c\}\). En efecto, si dos
palabras de peso seis con el mismo soporte no fueran proporcionales, se podría
escalar una de ellas para cancelar una coordenada al restarlas; se obtendría
una palabra no nula de \(C_W\) de peso menor que seis, contra la distancia
mínima. Como tampoco existe una palabra no nula igual a su opuesta en
característica tres, cada órbita proyectiva tiene exactamente dos elementos.
Por tanto,
\begin{equation}
 |\mathcal H_W|=\frac{264}{2}=132.
 \label{p12-83-c20-s2:eq:numero-hexadas-cw}
\end{equation}

Para \(T\in\binom{[12]}5\), definimos
\begin{equation}
 \nu(T):=\#\{H\in\mathcal H_W:T\subseteq H\}.
 \label{p12-83-c20-s2:eq:incidencia-cinco-hexada}
\end{equation}
La enumeración exacta de los
\[
 \binom{12}{5}=792
\]
subconjuntos de cinco puntos produce
\begin{equation}
 \nu(T)=1
 \qquad
 \text{para todo }T\in\binom{[12]}5.
 \label{p12-83-c20-s2:eq:unicidad-cinco-en-hexada}
\end{equation}

\begin{theorem}[Sistema de Witt ternario]
\label{p12-83-c20-s2:thm:w12-autonomo}
El sistema de incidencia
\[
 W_{12}:=([12],\mathcal H_W)
\]
es el sistema de Steiner
\begin{equation}
 \boxed{W_{12}=S(5,6,12).}
 \label{p12-83-c20-s2:eq:w12-steiner-autonomo}
\end{equation}
\end{theorem}

\begin{proof}
La igualdad \eqref{p12-83-c20-s2:eq:unicidad-cinco-en-hexada} es precisamente la
propiedad de Steiner. El número de bloques también se recupera de la doble
cuenta de pares \((T,H)\) con \(T\subset H\):
\[
 |\mathcal H_W|\binom65=\binom{12}{5},
 \qquad
 |\mathcal H_W|=\frac{792}{6}=132.
\]
\end{proof}

\subsection{Residuo de una tétrada y descomposición en \(4\) componentes incidentes}

Sea \(B\in\binom{[12]}4\). En un sistema \(S(5,6,12)\), el número de
hexadas que contienen \(B\) es
\begin{equation}
 \lambda_4
 =\frac{\binom{12-4}{5-4}}{\binom{6-4}{5-4}}
 =\frac82=4.
 \label{p12-83-c20-s2:eq:lambda-cuatro-w12}
\end{equation}
Escribimos estas cuatro hexadas como
\[
 H_i=B\sqcup P_i,
 \qquad |P_i|=2,
 \qquad i=1,\ldots,4.
\]

\begin{lemma}[Partición residual de una tétrada]
\label{p12-83-c20-s2:lem:particion-residual-tetrada-witt}
Los cuatro pares \(P_1,\ldots,P_4\) son disjuntos y
\begin{equation}
 [12]\setminus B
 =P_1\sqcup P_2\sqcup P_3\sqcup P_4.
 \label{p12-83-c20-s2:eq:particion-residual-tetrada-witt}
\end{equation}
\end{lemma}

\begin{proof}
Si dos pares residuales compartieran un punto \(p\), las correspondientes
hexadas contendrían el mismo subconjunto \(B\cup\{p\}\) de cinco puntos,
en contradicción con la unicidad de Steiner. Los cuatro pares son, pues,
disjuntos; contienen ocho puntos y agotan el complemento de \(B\).
\end{proof}

\begin{definition}[Estrella e involución de estrella de Witt]
\label{p12-83-c20-s2:def:estrella-witt-involucion}
La estrella centrada en \(B\) es
\begin{equation}
 \mathscr S_B
 :=\{H\in\mathcal H_W:B\subseteq H\}
 =\{B\sqcup P_1,\ldots,B\sqcup P_4\}.
 \label{p12-83-c20-s2:eq:estrella-witt-general}
\end{equation}
Si \(P_i=\{p_i^-,p_i^+\}\), su involución es
\begin{equation}
 \sigma_B:=\prod_{i=1}^{4}(p_i^-\ p_i^+).
 \label{p12-83-c20-s2:eq:involucion-estrella-witt-general}
\end{equation}
La permutación \(\sigma_B\) fija puntualmente \(B\) e intercambia los
extremos de cada pétalo residual.
\end{definition}

El término matemático empleado en adelante es \emph{estrella de Witt} o
\emph{involución de estrella de Witt}. No se trata de una estrella
euclídea.

\subsection{Estrella seleccionada por la tétrada del sello}

La unidad dodecafásica produjo, sin usar el diseño de Steiner,
\[
 B_K=\{4,6,9,10\}.
\]
Ahora que \(W_{12}\) está construido, su estrella queda bien definida. Las
cuatro hexadas incidentes son
\begin{align}
 H_{13}&=B_K\sqcup\{1,3\}
       =\{1,3,4,6,9,10\},
 \label{p12-83-c20-s2:eq:hexada-petalo-13}\\
 H_{2,11}&=B_K\sqcup\{2,11\}
       =\{2,4,6,9,10,11\},
 \label{p12-83-c20-s2:eq:hexada-petalo-211}\\
 H_{5,7}&=B_K\sqcup\{5,7\}
       =\{4,5,6,7,9,10\}
       =H_\alpha^-,
 \label{p12-83-c20-s2:eq:hexada-petalo-57}\\
 H_{8,12}&=B_K\sqcup\{8,12\}
       =\{4,6,8,9,10,12\}.
 \label{p12-83-c20-s2:eq:hexada-petalo-812}
\end{align}
En consecuencia,
\begin{equation}
 \boxed{
 \sigma_{B_K}=(1\ 3)(2\ 11)(5\ 7)(8\ 12).}
 \label{p12-83-c20-s2:eq:involucion-estrella-witt-bk}
\end{equation}

\begin{figure}[htbp]
\centering
\begin{tikzpicture}[
  >=Stealth,
  font=\small,
  core/.style={draw,double,rounded corners=2mm,align=center,
    minimum width=4.0cm,minimum height=1.3cm,fill=black!4},
  petal/.style={draw,rounded corners=2mm,align=center,
    minimum width=3.55cm,minimum height=1.08cm,fill=black!2},
  negative/.style={petal,line width=1.1pt,double,fill=black!10},
  arr/.style={->,thick}
]
\node[core] (B) at (0,0)
 {\(\displaystyle B_K=\{4,6,9,10\}\)\\
  puntos fijados por \(\sigma_{B_K}\)};
\node[petal] (P13) at (0,3.0)
 {\(P_1=\{1,3\}\), \((1\ 3)\)\\
  \(H_{13}=B_K\sqcup P_1\)};
\node[petal] (P211) at (5.25,0)
 {\(P_2=\{2,11\}\), \((2\ 11)\)\\
  \(H_{2,11}=B_K\sqcup P_2\)};
\node[negative] (P57) at (0,-3.0)
 {\(P_3=\{5,7\}\), \((5\ 7)\)\\
  \(H_{5,7}=H_\alpha^-\)};
\node[petal] (P812) at (-5.25,0)
 {\(P_4=\{8,12\}\), \((8\ 12)\)\\
  \(H_{8,12}=B_K\sqcup P_4\)};
\draw[arr] (B)--(P13);
\draw[arr] (B)--(P211);
\draw[arr] (B)--(P57);
\draw[arr] (B)--(P812);
\node[align=center,font=\footnotesize] at (0,-4.35)
 {\(\mathscr S_{B_K}
   =\{H_{13},H_{2,11},H_{5,7},H_{8,12}\}\).\\
  El diagrama representa incidencia, no distancia geométrica.};
\end{tikzpicture}
\caption{Estrella de Witt centrada en \(B_K\). El centro es la tétrada
fija; los cuatro pétalos son los pares residuales que determinan la
involución.}
\label{p12-83-c20-s2:fig:estrella-witt-bk-autonoma}
\end{figure}

\subsection{Campo global de las 495 estrellas}

Existe una estrella para cada tétrada, por lo que
\begin{equation}
 \#\{\mathscr S_B:B\in\tbinom{[12]}4\}
 =\binom{12}{4}=495.
 \label{p12-83-c20-s2:eq:numero-estrellas-witt}
\end{equation}
El censo exacto verifica
\begin{equation}
 \sigma_B(\mathcal H_W)=\mathcal H_W
 \qquad
 \text{para toda }B\in\binom{[12]}4.
 \label{p12-83-c20-s2:eq:estrellas-conservan-hexadas}
\end{equation}
Definimos
\begin{equation}
 \Gamma_\star
 :=\langle\sigma_B:B\in\tbinom{[12]}4\rangle
 \leq\operatorname{Aut}(W_{12}).
 \label{p12-83-c20-s2:eq:grupo-generado-estrellas-witt}
\end{equation}

Un conjunto generador de seis tétradas es
\[
 1234,
 \quad1235,
 \quad1236,
 \quad1245,
 \quad1345,
 \quad2345.
\]
Los órdenes de los subgrupos obtenidos al añadir sus involuciones
sucesivamente son
\begin{equation}
 1, 2, 6, 18, 360, 7920, 95040.
 \label{p12-83-c20-s2:eq:cadena-ordenes-generadores-m12}
\end{equation}

\begin{theorem}[Generación del grupo de Mathieu \(M_{12}\)]
\label{p12-83-c20-s2:thm:495-estrellas-m12-autonomo}
Se cumple
\begin{equation}
 \boxed{
 \langle\sigma_B:B\in\tbinom{[12]}4\rangle
 =\operatorname{Aut}(W_{12})
 \cong M_{12}.}
 \label{p12-83-c20-s2:eq:estrellas-generan-m12}
\end{equation}
\end{theorem}

\begin{proof}
Cada \(\sigma_B\) conserva las \(132\) hexadas, luego
\(\Gamma_\star\leq\operatorname{Aut}(W_{12})\). El censo de Schreier de
los seis generadores indicados produce orden \(95040\). El grupo de
automorfismos de \(S(5,6,12)\) es \(M_{12}\) y tiene ese mismo orden; por
tanto la inclusión es una igualdad.
\end{proof}

Una estrella aislada genera únicamente
\[
 \langle\sigma_B\rangle\cong C_2.
\]
La conclusión sobre \(M_{12}\) pertenece al campo global de las
\(495\) involuciones, no a \(\sigma_{B_K}\) considerada aisladamente.

\subsection{Acción sobre el módulo de aumentación}

Sea
\[
 V_{11}=\mathbf1^\perp\subset\mathbb R^{12}.
\]
La involución \(\sigma_B\) fija las cuatro coordenadas de \(B\) y actúa
por cuatro transposiciones sobre su complemento. En \(\mathbb R^{12}\)
posee ocho direcciones de autovalor \(+1\) y cuatro de autovalor \(-1\).
La dirección constante pertenece al primer autoespacio. Por tanto,
\begin{equation}
 \operatorname{Spec}(\sigma_B|_{V_{11}})
 =\{1^{[7]},(-1)^{[4]}\},
 \qquad
 \frac1{11}\operatorname{Tr}(\sigma_B|_{V_{11}})=\frac3{11}.
 \label{p12-83-c20-s2:eq:espectro-involucion-estrella-aumentacion}
\end{equation}
Es una involución firmada; no es un proyector positivo de rango tres.

\subsection{Obstrucciones de \(W_{12}\), del campo de \(495\) involuciones de estrella y de la generación de \(M_{12}\)}

\begin{criterion}[Refutación de \(W_{12}\), sus estrellas y \(M_{12}\)]
La construcción queda refutada si ocurre cualquiera de los hechos
siguientes:
\begin{enumerate}
 \item la proyectivización de las \(264\) palabras de peso seis no produce
       \(132\) soportes;
 \item existe un subconjunto de cinco puntos contenido en cero o en más de
       una hexada;
 \item una tétrada no pertenece exactamente a cuatro hexadas;
 \item los cuatro pares residuales no particionan el complemento de la
       tétrada;
 \item los pétalos de \(B_K\) no son
       \(\{1,3\},\{2,11\},\{5,7\},\{8,12\}\);
 \item alguna de las \(495\) involuciones no conserva
       \(\mathcal H_W\);
 \item el grupo generado no tiene orden \(95040\);
 \item se atribuye \(M_{12}\) a una única estrella en vez de al campo
       global de involuciones;
 \item se introduce la estrella antes de construir \(W_{12}\).
\end{enumerate}
\end{criterion}


% Unidad U020. Prolongacion binaria W12--W24 y M24.
% Redaccion autonoma desde la autoridad material 1063.

\section{Extensión binaria relativa a una dodecada, sistema \texorpdfstring{\(W_{24}\)}{W24} y grupo \texorpdfstring{\(M_{24}\)}{M24}}
\label{p12-83-c20-s3:chap:dodecada-w24-m24}

El sistema \(W_{12}\) construido sobre doce fases admite una prolongación
binaria cuando se declara un código de Golay extendido, una dodecada y un
isomorfismo de incidencia. Esos datos son entradas relativas de esta rama:
no se deducen únicamente de \(C_W\), y no intervienen en la rama reticular
ternaria que conducirá al retículo de Leech.

\subsection{Código binario de Golay y octadas}

Sea
\[
 \mathcal G_{24}\subset\mathbb F_2^{24}
\]
el código binario extendido de Golay. Sus pesos son
\begin{equation}
 0, 8, 12, 16, 24.
 \label{p12-83-c20-s3:eq:pesos-golay-binario-extendido}
\end{equation}
Los soportes de las palabras de peso ocho forman el sistema de Steiner
\begin{equation}
 W_{24}=S(5,8,24).
 \label{p12-83-c20-s3:eq:w24-steiner-octadas}
\end{equation}
Sus bloques se denominan \emph{octadas}. Fijemos una dodecada \(D\), es
decir, el soporte de una palabra de peso doce, y sea \(\mathcal O_{24}\)
la familia de octadas.

Definimos
\begin{equation}
 \mathcal H(D)
 :=\{O\cap D:O\in\mathcal O_{24},\ |O\cap D|=6\}.
 \label{p12-83-c20-s3:eq:hexadas-residuales-dodecada}
\end{equation}

\begin{theorem}[Diseño residual de una dodecada]
\label{p12-83-c20-s3:thm:diseno-residual-dodecada-autonomo}
El sistema residual es
\begin{equation}
 \boxed{(D,\mathcal H(D))\cong S(5,6,12).}
 \label{p12-83-c20-s3:eq:diseno-residual-dodecada-w12}
\end{equation}
Además, cada \(H\in\mathcal H(D)\) posee una única octada \(O_H\) tal que
\begin{equation}
 O_H\cap D=H.
 \label{p12-83-c20-s3:eq:elevacion-unica-hexada-octada}
\end{equation}
\end{theorem}

\begin{proof}
Sea \(A\subset D\), con \(|A|=5\). En \(W_{24}=S(5,8,24)\) existe una
única octada \(O_A\) que contiene \(A\). Si \(j=|O_A\cap D|\), la suma
binaria de las palabras características de \(O_A\) y \(D\) pertenece a
\(\mathcal G_{24}\) y tiene peso
\[
 8+12-2j=20-2j.
\]
Como \(A\subseteq O_A\cap D\), se tiene \(5\leq j\leq8\). Los cuatro
pesos posibles \(10,8,6,4\), combinados con
\eqref{p12-83-c20-s3:eq:pesos-golay-binario-extendido}, fuerzan \(j=6\). Así, todo
subconjunto de cinco puntos de \(D\) pertenece a una hexada residual única.

Si dos octadas tuvieran la misma intersección hexádica con \(D\), ambas
contendrían los mismos cinco puntos de esa hexada, contradiciendo la
unicidad en \(S(5,8,24)\). Esto prueba también la unicidad del
levantamiento.
\end{proof}

\subsection{Calibre relativo entre el diseño HMT y la dodecada}

Para relacionar la construcción anterior con el diseño obtenido de
\(C_W\), deben declararse explícitamente
\begin{equation}
 (D,\ell),
 \qquad
 \ell:W_{12}^{\rm HMT}
 \xrightarrow{\sim}(D,\mathcal H(D)),
 \label{p12-83-c20-s3:eq:dato-relativo-dodecada-alineamiento}
\end{equation}
donde \(\ell\) es un isomorfismo de incidencia. Si \(\ell_*\) denota su
acción sobre bloques, definimos
\begin{equation}
 \widehat\iota_{D,\ell}
 :=\iota_D\circ\ell_*:
 W_{12}^{\rm HMT}\longrightarrow W_{24},
 \qquad
 \widehat\iota_{D,\ell}(H)=O_{\ell(H)}.
 \label{p12-83-c20-s3:eq:elevacion-tipificada-w12-w24}
\end{equation}
La composición \(\widehat\iota_{D,\ell}\) evita aplicar
\(\iota_D\), cuyo dominio es \(\mathcal H(D)\), directamente a una hexada
que todavía vive en \([12]_{\rm HMT}\).

\begin{proposition}[Inyectividad de la elevación relativa]
\label{p12-83-c20-s3:prop:inyectividad-elevacion-w12-w24}
La aplicación \(\widehat\iota_{D,\ell}\) es inyectiva y conserva la
incidencia entre puntos y hexadas después de aplicar \(\ell\).
\end{proposition}

\begin{proof}
El isomorfismo \(\ell\) es biyectivo sobre bloques. Por el
\cref{p12-83-c20-s3:thm:diseno-residual-dodecada-autonomo}, cada hexada residual posee una
única octada elevadora. Dos hexadas HMT con la misma imagen tendrían la misma
hexada residual y, por biyectividad de \(\ell\), serían iguales.
\end{proof}

No existe una elección absoluta de \((D,\ell)\) previa a este calibre. El
objeto invariante es la clase de la construcción bajo cambio simultáneo de
dodecada y alineamiento.

\subsection{Ley de incidencia \texorpdfstring{\(4+1\)}{4+1} sobre una tétrada}

Para una tétrada \(B\subset D\), los parámetros derivados de los dos
sistemas de Steiner son
\begin{equation}
 \lambda_4(W_{12})
 =\frac{\binom81}{\binom21}=4,
 \qquad
 \lambda_4(W_{24})
 =\frac{\binom{20}1}{\binom41}=5.
 \label{p12-83-c20-s3:eq:ley-cuatro-mas-uno-witt}
\end{equation}
Por tanto, las cinco octadas que contienen \(B\) se descomponen de manera
única en
\begin{equation}
 \{O_H:H\in\mathcal H(D),\ B\subset H\}
 \sqcup\{O_B^\perp\}.
 \label{p12-83-c20-s3:eq:descomposicion-cinco-octadas}
\end{equation}
Las cuatro primeras satisfacen \(|O_H\cap D|=6\); la quinta satisface
\begin{equation}
 O_B^\perp\cap D=B.
 \label{p12-83-c20-s3:eq:octada-transversal-tetrada}
\end{equation}

Para \(B=B_K\), las cuatro elevaciones residuales corresponden a los
pétalos
\[
 \{1,3\},
 \qquad
 \{2,11\},
 \qquad
 \{5,7\},
 \qquad
 \{8,12\},
\]
y la quinta octada es la transversal \(O_{B_K}^\perp\). Esta ley explica
la transición de cuatro hexadas a cinco octadas sin identificar una cara
de un cubo con una hexada ni un vértice con una octada.

\subsection{Correspondencia de estratos y libertad del triple positivo}

La estructura de cinco regiones del canal \(\pi\) tiene tipos
\[
 1_\perp+1_{--}+3_{++}.
\]
La incidencia sobre \(B_K\) tiene tipos
\[
 1_{\rm transversal}+1_{H^-}+3_{\rm restantes}
\]
después de señalar la hexada negativa. Existe, por tanto, una
correspondencia natural entre \emph{estratos de rol}:
\begin{equation}
 1_\perp\longleftrightarrow O_{B_K}^\perp,
 \qquad
 1_{--}\longleftrightarrow
 \widehat\iota_{D,\ell}(H_\alpha^-),
 \qquad
 3_{++}\longleftrightarrow
 \mathcal O_{B_K}^{+,3}.
 \label{p12-83-c20-s3:eq:correspondencia-estratos-cinco-cinco}
\end{equation}
Aquí \(\mathcal O_{B_K}^{+,3}\) es el conjunto no ordenado de las otras
tres octadas residuales. Una biyección individual entre las tres regiones
positivas y esas tres octadas exige declarar un calibre adicional de
\(S_3\). La coincidencia del patrón \(1+1+3\) no determina por sí sola ese
ordenamiento.

\subsection{Grupo de automorfismos del sistema de octadas}

\begin{theorem}[Grupo de Mathieu \(M_{24}\)]
\label{p12-83-c20-s3:thm:automorfismos-w24-m24}
El grupo de automorfismos del sistema de Steiner de octadas es
\begin{equation}
 \boxed{
 \operatorname{Aut}(W_{24})\cong M_{24},
 \qquad
 |M_{24}|=244823040.}
 \label{p12-83-c20-s3:eq:automorfismos-w24-m24}
\end{equation}
\end{theorem}

El teorema identifica el grupo clásico de automorfismos del dato binario
de entrada. No afirma que \(M_{24}\) se genere a partir de una sola
dodecada, ni que el código ternario \(C_W\) determine por sí solo
\(\mathcal G_{24}\).

\subsection{Separación categórica respecto de la rama reticular}

La rama construida en esta unidad es
\begin{equation}
 (W_{12}^{\rm HMT};\mathcal G_{24},D,\ell)
 \longrightarrow W_{24}
 \longrightarrow\operatorname{Aut}(W_{24})\cong M_{24}.
 \label{p12-83-c20-s3:eq:rama-binaria-w24-m24}
\end{equation}
La rama reticular que se construirá después es
\begin{equation}
 C_W\longrightarrow N(A_2^{12})
 \xrightarrow{\text{vecino de índice tres}}\Lambda_{24}.
 \label{p12-83-c20-s3:eq:rama-ternaria-niemeier-leech-anunciada}
\end{equation}
Sus constructores, dominios y codominios son distintos. En particular,
\begin{equation}
 \boxed{
 \text{no se introduce ninguna flecha }
 W_{24}\longrightarrow\Lambda_{24}.}
 \label{p12-83-c20-s3:eq:no-flecha-w24-leech}
\end{equation}

\begin{figure}[htbp]
\centering
\begin{tikzpicture}[
  >=Stealth,
  node distance=13mm and 16mm,
  box/.style={draw,rounded corners=1.5mm,align=center,
    inner xsep=3mm,inner ysep=2.2mm},
  arr/.style={->,thick}
]
\node[box] (cw) {\(C_W\)};
\node[box,right=of cw] (w12) {\(W_{12}\)};
\node[box,right=of w12] (w24) {\(W_{24}\)};
\node[box,right=of w24] (m24) {\(M_{24}\)};
\node[box,below=of cw] (cw2) {\(C_W\)};
\node[box,right=of cw2] (n) {\(N(A_2^{12})\)};
\node[box,right=of n] (leech) {\(\Lambda_{24}\)};
\draw[arr] (cw)--(w12);
\draw[arr] (w12)--node[above,font=\scriptsize]
 {\((\mathcal G_{24},D,\ell)\)} (w24);
\draw[arr] (w24)--(m24);
\draw[arr] (cw2)--node[below,font=\scriptsize]
 {pegado ternario} (n);
\draw[arr] (n)--node[below,font=\scriptsize]
 {\(3\)-vecino} (leech);
\node[below=5mm of w24,font=\footnotesize,align=center]
 {La alineación vertical no es una flecha.};
\end{tikzpicture}
\caption{Bifurcación tipada de las prolongaciones binaria y ternaria.}
\label{p12-83-c20-s3:fig:bifurcacion-w24-leech-autonoma}
\end{figure}

\subsection{Obstrucciones de la prolongación binaria \(W_{12}\to W_{24}\): código de Golay, octadas, alineamiento y separación de la rama de Leech}

\begin{criterion}[Refutación de la prolongación binaria]
La construcción queda refutada si ocurre cualquiera de los hechos
siguientes:
\begin{enumerate}
 \item el espectro de pesos de \(\mathcal G_{24}\) contiene un peso no
       listado en \eqref{p12-83-c20-s3:eq:pesos-golay-binario-extendido};
 \item los soportes de peso ocho no forman \(S(5,8,24)\);
 \item una intersección residual de una octada con \(D\) que contiene cinco
       puntos de \(D\) tiene cardinal distinto de seis;
 \item una hexada residual admite cero o más de una octada elevadora;
 \item se aplica \(\iota_D\) directamente a una hexada HMT sin el
       alineamiento \(\ell\);
 \item la ley de incidencia sobre una tétrada no es \(4+1\);
 \item se ordena el triple positivo sin declarar el calibre \(S_3\);
 \item \(\operatorname{Aut}(W_{24})\) no tiene orden \(244823040\);
 \item se afirma que \(W_{24}\) surge de \(W_{12}\) sin el dato binario
       \(\mathcal G_{24}\), la dodecada y el alineamiento;
 \item se serializan las dos ramas mediante una flecha
       \(W_{24}\to\Lambda_{24}\).
\end{enumerate}
\end{criterion}


% Unidad U021. Pegado discriminante ternario y Niemeier A2^12.
% Redaccion autonoma desde la autoridad material 1063.

\section{Pegado discriminante ternario y construcción del retículo de Niemeier de tipo \texorpdfstring{\(A_2^{12}\)}{A₂¹²}}
\label{p12-83-c20-s4:chap:pegado-niemeier-a2-doce}

Esta unidad inicia la rama reticular directamente desde el código ternario
autodual \(C_W\). El sistema binario \(W_{24}\), la dodecada y el grupo
\(M_{24}\) no son entradas de esta construcción. La sucesión tipada es
\[
 C_W
 \longrightarrow
 (A_2^*/A_2)^{12}
 \longrightarrow
 N(C_W)
 \cong N(A_2^{12}).
\]

\subsection{Retículo raíz \texorpdfstring{\(A_2\)}{A₂}}

Sea
\[
 A_2=\mathbb Z\alpha_1\oplus\mathbb Z\alpha_2
\]
el retículo raíz cuya matriz de Gram, en la base ordenada
\((\alpha_1,\alpha_2)\), es
\begin{equation}
 G_{A_2}=
 \begin{pmatrix}
 2&-1\\
 -1&2
 \end{pmatrix}.
 \label{p12-83-c20-s4:eq:gram-a2-autonomo}
\end{equation}
Para \(x=x_1\alpha_1+x_2\alpha_2\),
\begin{equation}
 \langle x,x\rangle
 =2(x_1^2-x_1x_2+x_2^2).
 \label{p12-83-c20-s4:eq:norma-a2-autonoma}
\end{equation}
En particular, \(A_2\) es par, positivo definido y
\[
 \det A_2=3.
\]

Definimos
\[
 \rho:=\alpha_1+\alpha_2.
\]
Entonces
\begin{equation}
 \rho^2=2,
 \qquad
 \langle\rho,\alpha_1\rangle
 =\langle\rho,\alpha_2\rangle=1.
 \label{p12-83-c20-s4:eq:identidades-rho-a2}
\end{equation}

\subsection{Módulo y forma discriminantes}

El retículo dual es
\[
 A_2^*
 :=\{x\in A_2\otimes_{\mathbb Z}\mathbb Q:
       \langle x,A_2\rangle\subseteq\mathbb Z\}.
\]
Elegimos los representantes
\begin{equation}
 \mu_0=0,
 \qquad
 \mu_1=\frac{2\alpha_1+\alpha_2}{3},
 \qquad
 \mu_2=\frac{\alpha_1+2\alpha_2}{3}.
 \label{p12-83-c20-s4:eq:representantes-discriminantes-a2}
\end{equation}
Sus emparejamientos con la base simple son
\[
\begin{array}{c|cc}
 &\alpha_1&\alpha_2\\ \hline
 \mu_1&1&0\\
 \mu_2&0&1
\end{array},
\]
por lo que \(\mu_1,\mu_2\in A_2^*\). Además,
\[
 3\mu_1=2\alpha_1+\alpha_2\in A_2,
 \qquad
 3\mu_2=\alpha_1+2\alpha_2\in A_2,
\]
mientras \(\mu_1,\mu_2\notin A_2\). Se obtiene
\begin{equation}
 A_2^*/A_2
 =\{\overline\mu_0,\overline\mu_1,\overline\mu_2\}
 \cong\mathbb Z/3\mathbb Z
 \cong\mathbb F_3.
 \label{p12-83-c20-s4:eq:discriminante-a2-f3}
\end{equation}

La suma de representantes se reduce explícitamente mediante
\[
 \mu_1+\mu_1-\mu_2=\alpha_1,
 \qquad
 \mu_1+\mu_2=\rho,
 \qquad
 \mu_2+\mu_2-\mu_1=\alpha_2.
\]
Por consiguiente,
\begin{equation}
 \mu_r+\mu_s-\mu_{r+s\,({\rm mod}\,3)}\in A_2
 \qquad(r,s\in\mathbb F_3).
 \label{p12-83-c20-s4:eq:adicion-representantes-discriminantes-a2}
\end{equation}

\begin{proposition}[Forma discriminante de \(A_2\)]
\label{p12-83-c20-s4:prop:forma-discriminante-a2}
Bajo \eqref{p12-83-c20-s4:eq:discriminante-a2-f3}, las formas bilineal y cuadrática
discriminantes son
\begin{align}
 b(\overline\mu_r,\overline\mu_s)
 &\equiv\frac{2rs}{3}\pmod{\mathbb Z},
 \label{p12-83-c20-s4:eq:bilineal-discriminante-a2}\\
 q(\overline\mu_r)
 &\equiv\frac{2r^2}{3}\pmod{2\mathbb Z}.
 \label{p12-83-c20-s4:eq:cuadratica-discriminante-a2}
\end{align}
\end{proposition}

\begin{proof}
La fórmula \eqref{p12-83-c20-s4:eq:norma-a2-autonoma} da
\[
 \mu_1^2=\mu_2^2=\frac23,
 \qquad
 \langle\mu_1,\mu_2\rangle=\frac13.
\]
Para \(r,s\in\{0,1,2\}\), estos valores coinciden, módulo
\(\mathbb Z\), con \(2rs/3\). La fórmula cuadrática resulta al tomar
\(r=s\) y recordar que el retículo es par.
\end{proof}

\subsection{\(12\) factores y aplicación de pegado}

Sea
\[
 L_0:=A_2^{12}.
\]
Entonces
\[
 L_0^*=(A_2^*)^{12},
 \qquad
 L_0^*/L_0\cong\mathbb F_3^{12}.
\]
Para \(c=(c_1,\ldots,c_{12})\in\mathbb F_3^{12}\), definimos
\begin{equation}
 \phi(c):=(\mu_{c_1},\ldots,\mu_{c_{12}})\in L_0^*.
 \label{p12-83-c20-s4:eq:aplicacion-pegado-cw}
\end{equation}
La extensión determinada por \(C_W\) es
\begin{equation}
 \begin{aligned}
 N(C_W)
 &:=\bigcup_{c\in C_W}(L_0+\phi(c))\\
 &=\{x\in L_0^*:x\bmod L_0\in C_W\}.
 \end{aligned}
 \label{p12-83-c20-s4:eq:definicion-niemeier-cw}
\end{equation}

\begin{lemma}[Clausura aditiva del pegado]
\label{p12-83-c20-s4:lem:clausura-pegado-cw}
El conjunto \(N(C_W)\) es un subgrupo discreto de \(L_0^*\) de rango
veinticuatro.
\end{lemma}

\begin{proof}
Sean \(x=\ell+\phi(c)\) y \(y=m+\phi(d)\), con \(\ell,m\in L_0\) y
\(c,d\in C_W\). Por \eqref{p12-83-c20-s4:eq:adicion-representantes-discriminantes-a2},
\[
 \phi(c)+\phi(d)-\phi(c+d)\in L_0.
\]
Como \(C_W\) es lineal, \(c+d\in C_W\), y entonces
\(x+y\in L_0+\phi(c+d)\). La misma identidad aplicada al opuesto da
inversos. Las inclusiones
\[
 L_0\subseteq N(C_W)\subseteq L_0^*
\]
prueban discreción y rango veinticuatro.
\end{proof}

\subsection{Integralidad inducida por autodualidad}

La forma bilineal discriminante de \(L_0\) es la suma ortogonal de doce
copias de \eqref{p12-83-c20-s4:eq:bilineal-discriminante-a2}. Para
\(c,d\in\mathbb F_3^{12}\),
\begin{equation}
 b(c,d)
 \equiv\frac23\sum_{i=1}^{12}c_id_i
 \pmod{\mathbb Z}.
 \label{p12-83-c20-s4:eq:bilineal-discriminante-a2-doce}
\end{equation}

\begin{lemma}[Integralidad del pegado]
\label{p12-83-c20-s4:lem:integralidad-pegado-cw}
El retículo \(N(C_W)\) es integral.
\end{lemma}

\begin{proof}
La autodualidad \(C_W=C_W^\perp\) implica
\[
 \sum_{i=1}^{12}c_id_i\equiv0\pmod3
 \qquad(c,d\in C_W).
\]
Por \eqref{p12-83-c20-s4:eq:bilineal-discriminante-a2-doce}, la forma discriminante se
anula sobre \(C_W\times C_W\). Por tanto, el emparejamiento de cualesquiera
dos clases pegadas es entero.
\end{proof}

\subsection{Paridad inducida por los pesos del código}

En una coordenada no nula del código, la forma cuadrática discriminante
vale \(2/3\) módulo \(2\mathbb Z\). Por tanto,
\begin{equation}
 q(c)\equiv\frac23\operatorname{wt}(c)\pmod{2\mathbb Z}.
 \label{p12-83-c20-s4:eq:cuadratica-discriminante-cw}
\end{equation}

\begin{lemma}[Paridad del pegado]
\label{p12-83-c20-s4:lem:paridad-pegado-cw}
El retículo \(N(C_W)\) es par.
\end{lemma}

\begin{proof}
Los únicos pesos de \(C_W\) son \(0,6,9,12\), y
\[
 \frac23\{0,6,9,12\}=\{0,4,6,8\}\subseteq2\mathbb Z.
\]
Toda clase de pegado es isotrópica para la forma cuadrática discriminante,
de modo que todas las normas de \(N(C_W)\) son pares.
\end{proof}

Sean \(g_1,\ldots,g_6\) las filas de \(G_W=(I_6\mid A_W)\) y
\(\gamma_i=\phi(g_i)\). La multiplicación exacta en \(A_2^*\) da
\begin{equation}
 (\langle\gamma_i,\gamma_j\rangle)_{i,j=1}^{6}
 =
 \begin{pmatrix}
 4&2&2&2&2&2\\
 2&4&2&2&2&2\\
 2&2&4&2&2&2\\
 2&2&2&4&2&2\\
 2&2&2&2&4&2\\
 2&2&2&2&2&4
 \end{pmatrix}.
 \label{p12-83-c20-s4:eq:gram-generadores-pegado-cw}
\end{equation}
Esta matriz proporciona un control directo adicional de integralidad y
paridad para los seis generadores de pegado.

\subsection{Índice, determinante y unimodularidad}

Las clases laterales de \(L_0\) que aparecen en
\eqref{p12-83-c20-s4:eq:definicion-niemeier-cw} están indexadas por los
\(|C_W|=729=3^6\) elementos del código y son distintas. En consecuencia,
\begin{equation}
 [N(C_W):L_0]=3^6.
 \label{p12-83-c20-s4:eq:indice-niemeier-sobre-a2-doce}
\end{equation}
Como
\[
 \det L_0=(\det A_2)^{12}=3^{12},
\]
la fórmula del determinante para una extensión finita da
\begin{equation}
 \det N(C_W)
 =\frac{\det L_0}{[N(C_W):L_0]^2}
 =\frac{3^{12}}{(3^6)^2}=1.
 \label{p12-83-c20-s4:eq:determinante-niemeier-cw}
\end{equation}

\subsection{Sistema de raíces inducido por la incidencia de Witt y clasificación de sus órbitas}

En cada una de las dos clases discriminantes no nulas de
\(A_2^*/A_2\), la norma mínima es
\[
 \mu_1^2=\mu_2^2=\frac23.
\]
Si \(c\in C_W\setminus\{0\}\), entonces
\(\operatorname{wt}(c)\geq6\). Por ortogonalidad de los doce factores,
cualquier vector de \(L_0+\phi(c)\) tiene norma al menos
\begin{equation}
 6\cdot\frac23=4.
 \label{p12-83-c20-s4:eq:cota-clases-no-triviales-niemeier}
\end{equation}
Por tanto, ninguna clase de pegado no trivial contiene vectores de norma
dos. Las raíces de \(N(C_W)\) son exactamente las raíces ya presentes en
\(L_0=A_2^{12}\).

\begin{theorem}[Realización del Niemeier de tipo \(A_2^{12}\)]
\label{p12-83-c20-s4:thm:niemeier-a2-doce-autonomo}
El retículo \(N(C_W)\) es positivo definido, par y unimodular, y su sistema
de raíces es \(A_2^{12}\). En consecuencia,
\begin{equation}
 \boxed{N(C_W)\cong N(A_2^{12}).}
 \label{p12-83-c20-s4:eq:identificacion-niemeier-a2-doce}
\end{equation}
\end{theorem}

\begin{proof}
La positividad procede de \(L_0^*\). La
\cref{p12-83-c20-s4:lem:integralidad-pegado-cw,p12-83-c20-s4:lem:paridad-pegado-cw} demuestra
integralidad y paridad; \eqref{p12-83-c20-s4:eq:determinante-niemeier-cw} demuestra
unimodularidad. La cota \eqref{p12-83-c20-s4:eq:cota-clases-no-triviales-niemeier}
identifica el sistema de raíces. La clasificación de Niemeier asocia de
manera única a ese sistema de raíces el retículo indicado.
\end{proof}

\subsection{Obstrucciones del pegado ternario \(C_W\to N(A_2^{12})\): paridad, unimodularidad y sistema de raíces}

\begin{criterion}[Refutación del pegado ternario]
La construcción queda refutada si ocurre cualquiera de los hechos
siguientes:
\begin{enumerate}
 \item alguno de los representantes \(\mu_0,\mu_1,\mu_2\) no pertenece a
       \(A_2^*\), dos de sus clases coinciden, o falla
       \eqref{p12-83-c20-s4:eq:adicion-representantes-discriminantes-a2};
 \item la forma discriminante no es la dada en
       \cref{p12-83-c20-s4:prop:forma-discriminante-a2};
 \item \(N(C_W)\) no es cerrado bajo suma o inversos;
 \item existe \(c,d\in C_W\) con
       \(\frac23\sum_i c_id_i\notin\mathbb Z\);
 \item una palabra de \(C_W\) produce
       \(\frac23\operatorname{wt}(c)\notin2\mathbb Z\);
 \item el índice no es \(3^6\) o el determinante no es uno;
 \item una clase de pegado no trivial contiene un vector de norma dos;
 \item se utiliza \(W_{24}\), una dodecada o un dato binario en cualquiera
       de las operaciones de pegado anteriores.
\end{enumerate}
\end{criterion}


% Unidad U022. Origen incidencial, vecino ternario y Leech.
% Redaccion autonoma desde la autoridad material 1063.

\section{Origen incidencial, problema radial tipado y construcción del vecino unimodular sin raíces}
\label{p12-83-c20-s5:chap:origen-incidencial-vecino-leech}

El retículo \(N(A_2^{12})\) ya está construido. Esta unidad determina un
punto base a partir de un marco ordenado de incidencias, resuelve un
problema variacional en un dominio declarado, verifica las tres condiciones
del vecino de índice tres y demuestra que el resultado es el retículo de
Leech. El vector del vecino se denota \(v_{\mathcal F_*}\): depende del
marco incidencial y no interviene en el acarreo numérico de \(\alpha\).

\subsection{Marco ordenado de incidencias}

Consideremos las cuatro hexadas
\begin{align}
 H^-&=\{4,5,6,7,9,10\},
 \label{p12-83-c20-s5:eq:hexada-marco-negativa}\\
 H_e&=\{1,3,4,6,9,10\},
 \label{p12-83-c20-s5:eq:hexada-marco-e}\\
 H_\varphi&=\{1,2,3,4,7,9\},
 \label{p12-83-c20-s5:eq:hexada-marco-phi}\\
 H_{\rm APP}&=\{2,3,5,10,11,12\}.
 \label{p12-83-c20-s5:eq:hexada-marco-app}
\end{align}
El dato es el marco ordenado
\begin{equation}
 \mathcal F:=(H^-,H_e,H_\varphi,H_{\rm APP}).
 \label{p12-83-c20-s5:eq:marco-incidencial-ordenado}
\end{equation}
Definimos
\begin{equation}
 o(\mathcal F)
 :=(H_e\cap H_\varphi)\setminus(H^-\cup H_{\rm APP}).
 \label{p12-83-c20-s5:eq:operador-origen-incidencial}
\end{equation}

\begin{theorem}[Reconstrucción equivariante del origen]
\label{p12-83-c20-s5:thm:origen-incidencial-equivariante}
Para el marco canónico \(\mathcal F_*\),
\begin{equation}
 \boxed{o(\mathcal F_*)=\{1\}.}
 \label{p12-83-c20-s5:eq:origen-incidencial-singleton}
\end{equation}
Además, para toda permutación simultánea \(g\in M_{12}\),
\begin{equation}
 o(g\mathcal F)=g\,o(\mathcal F).
 \label{p12-83-c20-s5:eq:equivarianza-origen-incidencial}
\end{equation}
\end{theorem}

\begin{proof}
Se tiene
\[
 H_e\cap H_\varphi=\{1,3,4,9\}
\]
y
\[
 (H^-\cup H_{\rm APP})\cap\{1,3,4,9\}=\{3,4,9\}.
\]
La diferencia es \(\{1\}\). Para la equivarianza se aplican las
identidades
\[
 g(A\cap B)=gA\cap gB,
 \quad
 g(A\cup B)=gA\cup gB,
 \quad
 g(A\setminus B)=gA\setminus gB,
\]
válidas para toda biyección.
\end{proof}

La tétrada \(B_K=\{4,6,9,10\}\) no determina por sí sola el origen. El
singleton se obtiene del marco ordenado completo; no es una coordenada
elegida después de observar el resultado.

\subsection{Cámara radial positiva del sistema de raíces y selección del cono dominante}

Sea \(\rho_i\) la copia de \(\rho=\alpha_1+\alpha_2\) en el
\(i\)-ésimo factor de \(A_2^{12}\). Consideramos el dominio
\begin{equation}
 \mathscr R_+
 :=\left\{
 v(a)=\sum_{i=1}^{12}a_i\rho_i:
 a_i=1+3b_i, b_i\in\mathbb Z_{\geq0}
 \right\}.
 \label{p12-83-c20-s5:eq:camara-radial-positiva}
\end{equation}
Como los factores son ortogonales y \(\rho_i^2=2\), si
\(S=\sum_i b_i\), entonces
\begin{align}
 v(a)^2
 &=2\sum_{i=1}^{12}(1+3b_i)^2\notag\\
 &=24+12S+18\sum_{i=1}^{12}b_i^2.
 \label{p12-83-c20-s5:eq:norma-radial-expandida}
\end{align}
Reduciendo módulo dieciocho,
\[
 v(a)^2\equiv6+12S\pmod{18},
\]
y por tanto
\begin{equation}
 v(a)^2\equiv0\pmod{18}
 \quad\Longleftrightarrow\quad
 S\equiv1\pmod3.
 \label{p12-83-c20-s5:eq:congruencia-radial-suma-b}
\end{equation}

\begin{proposition}[Mínimos en la cámara radial declarada]
\label{p12-83-c20-s5:prop:minimos-camara-radial}
En el subconjunto de \(\mathscr R_+\) definido por
\(v(a)^2\equiv0\pmod{18}\), la norma mínima es \(54\). Se alcanza
exactamente en doce vectores: para cada \(j\in\{1,\ldots,12\}\), el
coeficiente \(a_j\) vale cuatro y los once restantes valen uno.
\end{proposition}

\begin{proof}
La condición \eqref{p12-83-c20-s5:eq:congruencia-radial-suma-b} y la no negatividad de
los \(b_i\) implican que el menor valor posible de \(S\) es uno. Esto
obliga a un único índice \(j\) con \(b_j=1\) y a \(b_i=0\) para
\(i\neq j\). Sustituyendo en \eqref{p12-83-c20-s5:eq:norma-radial-expandida},
\[
 v(a)^2=24+12+18=54.
\]
Hay doce elecciones de \(j\). El siguiente valor admisible de \(S\) es
cuatro y produce una norma estrictamente mayor.
\end{proof}

El origen \(o(\mathcal F_*)=\{1\}\) selecciona
\begin{equation}
 \boxed{
 v_{\mathcal F_*}
 =4\rho_1+\rho_2+\cdots+\rho_{12},
 \qquad
 v_{\mathcal F_*}^2=54.}
 \label{p12-83-c20-s5:eq:vector-marco-incidencial}
\end{equation}

\begin{remark}[Alcance exacto de la minimalidad]
La norma \(54\) es mínima en la cámara radial positiva
\eqref{p12-83-c20-s5:eq:camara-radial-positiva}, no en toda la clase residual. El vector
\[
 u=(\alpha_1-2\alpha_2,\rho,\ldots,\rho)
\]
satisface
\[
 \alpha_1-2\alpha_2-\rho=-3\alpha_2,
\]
de modo que
\[
 u\equiv(\rho,\ldots,\rho)
 \equiv v_{\mathcal F_*}\pmod{3A_2^{12}}.
\]
Sin embargo,
\[
 (\alpha_1-2\alpha_2)^2
 =2(1+2+4)=14,
 \qquad
 u^2=14+11\cdot2=36.
\]
Este testigo impide extender la afirmación de minimalidad fuera del dominio
de \cref{p12-83-c20-s5:prop:minimos-camara-radial}.
\end{remark}

\subsection{Carácter y núcleo de congruencia}

Escribimos
\[
 N:=N(A_2^{12}),
 \qquad
 v:=v_{\mathcal F_*}.
\]
Definimos
\begin{equation}
 \chi_v:N\longrightarrow\mathbb Z/3\mathbb Z,
 \qquad
 \chi_v(x)=\langle x,v\rangle\bmod3,
 \label{p12-83-c20-s5:eq:caracter-vecino-leech}
\end{equation}
su núcleo
\begin{equation}
 N_0:=\ker\chi_v
 =\{x\in N:\langle x,v\rangle\equiv0\pmod3\},
 \label{p12-83-c20-s5:eq:nucleo-congruencia-vecino}
\end{equation}
y la extensión
\begin{equation}
 N_v:=N_0+\mathbb Z\frac v3.
 \label{p12-83-c20-s5:eq:definicion-vecino-indice-tres}
\end{equation}

\subsection{Condición I1: sobreyectividad y exclusión de raíces}

Por construcción, \(v\in A_2^{12}\subseteq N\). Las raíces de un factor
\(A_2\) son
\[
 \pm\alpha_1,
 \quad
 \pm\alpha_2,
 \quad
 \pm\rho.
\]
Sus emparejamientos con \(\rho\) son \(\pm1\) o \(\pm2\). Todos los
coeficientes de \(v\) son congruentes con uno módulo tres:
\[
 4\equiv1\pmod3,
 \qquad
 1\equiv1\pmod3.
\]
Por tanto, ninguna raíz de \(A_2^{12}\) tiene emparejamiento nulo con
\(v\) módulo tres. En particular,
\[
 \chi_v(\alpha_{1,1})=1,
\]
de modo que \(\chi_v\) es sobreyectivo y
\begin{equation}
 [N:N_0]=3.
 \label{p12-83-c20-s5:eq:indice-nucleo-vecino}
\end{equation}
Además, ninguna raíz de \(N\) pertenece a \(N_0\).

\subsection{Condición I2: divisibilidad de la norma}

La identidad \eqref{p12-83-c20-s5:eq:vector-marco-incidencial} da
\begin{equation}
 v^2=54\equiv0\pmod{18},
 \qquad
 \left(\frac v3\right)^2=6.
 \label{p12-83-c20-s5:eq:condicion-i2-vecino}
\end{equation}
En particular, \(\langle v,v\rangle\equiv0\pmod3\), luego
\(v\in N_0\).

\subsection{Condición I3: clase dual de orden \(3\)}

Si \(x\in N_0\), existe \(k\in\mathbb Z\) con
\(\langle x,v\rangle=3k\). Por tanto,
\[
 \left\langle x,\frac v3\right\rangle=k\in\mathbb Z,
\]
y
\begin{equation}
 \frac v3\in N_0^*.
 \label{p12-83-c20-s5:eq:v-tercio-en-dual-nucleo}
\end{equation}
Por otra parte,
\[
 \left\langle\frac v3,\alpha_{1,1}\right\rangle=\frac43.
\]
Como \(\alpha_{1,1}\in N\) y \(N\) es integral, esto demuestra
\begin{equation}
 \frac v3\notin N.
 \label{p12-83-c20-s5:eq:v-tercio-no-en-niemeier}
\end{equation}
Finalmente, \(3(v/3)=v\in N_0\). La clase de \(v/3\) en
\(N_0^*/N_0\) es no trivial y tiene orden exactamente tres. En
consecuencia,
\begin{equation}
 [N_v:N_0]=3.
 \label{p12-83-c20-s5:eq:indice-extension-vecino}
\end{equation}

\subsection{Paridad, integralidad y determinante del vecino}

\begin{proposition}[Estructura unimodular del vecino]
\label{p12-83-c20-s5:prop:estructura-unimodular-vecino}
El retículo \(N_v\) es integral, par y unimodular.
\end{proposition}

\begin{proof}
Como \(N\) es unimodular y \([N:N_0]=3\),
\[
 \det N_0=[N:N_0]^2\det N=9.
\]
Como \([N_v:N_0]=3\),
\[
 \det N_v=\frac{\det N_0}{[N_v:N_0]^2}=1.
\]

Sea \(y=x+m(v/3)\), con \(x\in N_0\), y escribamos
\(\langle x,v\rangle=3k\). Entonces
\begin{align*}
 y^2
 &=x^2+2m\left\langle x,\frac v3\right\rangle
       +m^2\left(\frac v3\right)^2\\
 &=x^2+2mk+6m^2.
\end{align*}
Los tres sumandos forman un entero par. La forma polarizada prueba
integralidad, y el determinante uno prueba unimodularidad.
\end{proof}

\subsection{\(2\) clases globales y \(6\) clases locales}

Globalmente, \eqref{p12-83-c20-s5:eq:definicion-vecino-indice-tres} añade sólo dos clases
laterales no triviales:
\[
 N_0+\frac v3,
 \qquad
 N_0-\frac v3.
\]
Localmente, en cada factor \(A_2\), una coordenada de esas clases pertenece
a una de las seis clases
\begin{equation}
 \mathscr C_{j,\pm}
 :=A_2+\mu_j\pm\frac\rho3,
 \qquad j\in\{0,1,2\}.
 \label{p12-83-c20-s5:eq:seis-clases-locales-vecino}
\end{equation}
En la primera coordenada, \(4\rho/3\equiv\rho/3\pmod{A_2}\), por lo que
la misma lista vale para los doce factores.

\begin{table}[htbp]
\centering
\caption{Representantes mínimos de las seis clases laterales locales.}
\label{p12-83-c20-s5:tab:seis-clases-locales-vecino}
\begin{tabular}{cclc}
\toprule
clase&representante mínimo&\((u,v)\)&
\(2(u^2-uv+v^2)/9\)\\
\midrule
\(\mathscr C_{0,-}\)&\(-\rho/3\)&\((-1,-1)\)&\(2/9\)\\
\(\mathscr C_{0,+}\)&\( \rho/3\)&\((1,1)\)&\(2/9\)\\
\(\mathscr C_{1,-}\)&\( \mu_1-\rho/3\)&\((1,0)\)&\(2/9\)\\
\(\mathscr C_{1,+}\)&\( \mu_1+\rho/3-\rho\)&\((0,-1)\)&\(2/9\)\\
\(\mathscr C_{2,-}\)&\( \mu_2-\rho/3\)&\((0,1)\)&\(2/9\)\\
\(\mathscr C_{2,+}\)&\( \mu_2+\rho/3-\rho\)&\((-1,0)\)&\(2/9\)\\
\bottomrule
\end{tabular}
\end{table}

Aquí \((u,v)\) significa el representante
\((u\alpha_1+v\alpha_2)/3\).

\begin{lemma}[Mínimo exacto de las seis clases]
\label{p12-83-c20-s5:lem:minimo-seis-clases-locales}
Cada clase de \eqref{p12-83-c20-s5:eq:seis-clases-locales-vecino} tiene norma mínima
exactamente \(2/9\).
\end{lemma}

\begin{proof}
La norma de \((u\alpha_1+v\alpha_2)/3\) es
\[
 \frac29(u^2-uv+v^2).
\]
La forma entera \(u^2-uv+v^2\) es positiva definida. En una clase residual
no nula no toma el valor cero, luego es al menos uno. Los seis
representantes de la \cref{p12-83-c20-s5:tab:seis-clases-locales-vecino} alcanzan uno.
\end{proof}

\subsection{Criterio de ausencia de raíces en el retículo de Leech y consecuencia para la distancia mínima}

\begin{theorem}[Ausencia de vectores de norma dos]
\label{p12-83-c20-s5:thm:vecino-sin-raices}
El retículo \(N_v\) no contiene vectores de norma dos.
\end{theorem}

\begin{proof}
Todas las raíces de \(N\) son las de \(A_2^{12}\). La condición I1 prueba
que ninguna pertenece a \(N_0\). Por tanto, la clase \(N_0\) no contiene
vectores de norma dos.

En una clase nueva \(N_0\pm v/3\), cada una de las doce coordenadas
pertenece a una de las seis clases locales. Por
\cref{p12-83-c20-s5:lem:minimo-seis-clases-locales}, cada coordenada contribuye al menos
\(2/9\). La ortogonalidad da
\[
 \left\|x\pm\frac v3\right\|^2
 \geq12\cdot\frac29=\frac83>2.
\]
Tampoco las clases nuevas contienen raíces.
\end{proof}

\subsection{Identificación con el retículo de Leech}

\begin{theorem}[Construcción del retículo de Leech]
\label{p12-83-c20-s5:thm:identificacion-vecino-leech}
El retículo
\begin{equation}
 N_v
 =\ker(\langle\,\cdot\,,v_{\mathcal F_*}\rangle\bmod3)
  +\mathbb Z\frac{v_{\mathcal F_*}}3
 \label{p12-83-c20-s5:eq:formula-final-vecino-leech}
\end{equation}
es isométrico al retículo de Leech:
\begin{equation}
 \boxed{N_v\cong\Lambda_{24}.}
 \label{p12-83-c20-s5:eq:identificacion-leech-autonoma}
\end{equation}
\end{theorem}

\begin{proof}
El retículo \(N_v\) tiene rango veinticuatro y es positivo definido. La
\cref{p12-83-c20-s5:prop:estructura-unimodular-vecino} demuestra que es par y unimodular;
el \cref{p12-83-c20-s5:thm:vecino-sin-raices} demuestra que carece de raíces. Por la
clasificación de Niemeier, existe una única clase de isometría de retículos
positivos definidos, pares y unimodulares de rango veinticuatro sin raíces:
el retículo de Leech.
\end{proof}

\subsection{Los 243 síndromes del código ternario perforado}

Sea \(G_{11}\) el código ternario de Golay de parámetros
\([11,6,5]_3\). Su espacio de síndromes tiene cardinal
\begin{equation}
 3^{11-6}=3^5=243.
 \label{p12-83-c20-s5:eq:numero-sindromes-golay-ternario}
\end{equation}
El número de errores de peso a lo sumo dos es
\begin{equation}
 1+2\binom{11}{1}+2^2\binom{11}{2}
 =1+22+220=243.
 \label{p12-83-c20-s5:eq:bola-radio-dos-golay-ternario}
\end{equation}

\begin{proposition}[Correspondencia perfecta error--síndrome]
\label{p12-83-c20-s5:prop:correspondencia-error-sindrome-golay}
Cada síndrome de \(G_{11}\) posee un único representante de error de peso
a lo sumo dos.
\end{proposition}

\begin{proof}
La distancia mínima cinco implica que dos errores distintos de peso a lo
sumo dos no pueden diferir en una palabra del código: su diferencia tendría
peso a lo sumo cuatro. La aplicación error--síndrome es, por tanto,
inyectiva sobre la bola de radio dos. Los dos conjuntos tienen cardinal
\(243\) por \eqref{p12-83-c20-s5:eq:numero-sindromes-golay-ternario} y
\eqref{p12-83-c20-s5:eq:bola-radio-dos-golay-ternario}; la aplicación es biyectiva.
\end{proof}

Estos \(243\) síndromes no son las \(243\) palabras visibles obtenidas del
catálogo TPK. Comparten un cardinal, pero poseen dominios, equivalencias y
funciones distintas.

\subsection{Obstrucciones del vecino ternario \(N_v\cong\Lambda_{24}\): clase de orden \(3\), ausencia de raíces y separación de los cardinales \(243\)}

\begin{criterion}[Refutación del vecino y de la identificación de Leech]
La construcción queda refutada si ocurre cualquiera de los hechos
siguientes:
\begin{enumerate}
 \item el operador \(o(\mathcal F)\) no produce \(\{1\}\) o no es
       equivariante;
 \item existe en la cámara radial un vector admisible de norma menor que
       \(54\), o el número de minimizadores no es doce;
 \item se afirma minimalidad global pese al testigo de norma \(36\);
 \item \(\chi_v\) no es sobreyectivo o una raíz pertenece a \(N_0\);
 \item \(v^2\not\equiv0\pmod{18}\) o \((v/3)^2\neq6\);
 \item \(v/3\notin N_0^*\), \(v/3\in N\), o su clase no tiene orden tres;
 \item el determinante de \(N_v\) no es uno o el vecino deja de ser par;
 \item alguna de las seis clases locales tiene mínimo menor que \(2/9\);
 \item una de las dos clases globales nuevas contiene una raíz;
 \item se confunden las dos clases globales con las seis clases locales;
 \item se utiliza \(W_{24}\) como antecedente del pegado o del vecino;
 \item los \(243\) síndromes se identifican con las \(243\) palabras
       visibles por compartir cardinal.
\end{enumerate}
\end{criterion}


\section{Holonomía finita, configuración de Schläfli y simetría \texorpdfstring{\(E_6\)}{E6}}
\label{p12-83-c20-s6:chap:holonomia-schlafli-e6-nuevo}
\index{holonomía finita}
\index{Schläfli}
\index{E6@\(E_6\)}

La rama de Witt no agota las estructuras excepcionales accesibles desde
APP. La rama paralela de este capítulo parte de las dos polarizaciones
mixtas suma--producto y producto--suma, junto con una orientación de
plaqueta y una sección central declaradas. La curvatura calculada determina
la forma alternante; ésta define la extensión central. No se infiere una
extensión a partir de la sola cardinalidad 27.

\subsection{Curvaturas mixtas de APP}

Sobre \(T_9=(\mathbb Z/9\mathbb Z)^2\), con
\(S(x,y)=x+y\) y \(P(x,y)=xy\), sean
\[
 \Delta_xf=f(x+1,y)-f(x,y),\qquad
 \Delta_yf=f(x,y+1)-f(x,y).
\]
Las conexiones mixtas son
\[
 A^{SP}=(\Delta_xS,\Delta_yP)=(1,x),\qquad
 A^{PS}=(\Delta_xP,\Delta_yS)=(y,1).
\]
Su curvatura discreta es
\[
 F_A=\Delta_xA_y-\Delta_yA_x.
\]

\begin{proposition}[Curvaturas orientadas opuestas]
\label{p12-83-c20-s6:prop:curvaturas-app-e6-nuevo}
\[
 F_{A^{SP}}=1,\qquad F_{A^{PS}}=-1\pmod9.
\]
La circulación sobre un rectángulo orientado de lados \(r,s\) es,
respectivamente, \(rs\) y \(-rs\).
\end{proposition}

\begin{proof}
\(\Delta_xx-\Delta_y1=1\) y
\(\Delta_x1-\Delta_yy=-1\). Un rectángulo contiene \(rs\) plaquetas; al
sumar, las aristas interiores se cancelan y queda la circulación del borde.
\end{proof}

Para desplazamientos \(u=(a,b)\), \(v=(c,d)\), la clase alternante es
\begin{equation}
 \sigma(u,v)=ad-bc\pmod9.
 \label{p12-83-c20-s6:eq:forma-alternante-e6-nueva}
\end{equation}
La orientación fija el signo de \(\sigma\). La sección central empleada
abajo pertenece al calibre de esta realización; cambiarla por un coborde
produce una extensión isomorfa, no un nuevo invariante.

\subsection{Extensión central y reducción ternaria}

Definimos el grupo de Heisenberg nonádico
\begin{equation}
 \mathcal H_9=\{(a,b,z):a,b,z\in\mathbb Z/9\mathbb Z\},
 \quad
 (a,b,z)(c,d,w)=(a+c,b+d,z+w+bc).
 \label{p12-83-c20-s6:eq:heisenberg-9-nuevo}
\end{equation}
El conmutador de dos desplazamientos horizontales es
\[
 [(a,b,0),(c,d,0)]=(0,0,bc-ad),
\]
que registra \(-\sigma(u,v)\). Reduciendo cada coordenada módulo tres se
obtiene
\begin{equation}
 H_3=\{(a,b,z):a,b,z\in\mathbb F_3\},\qquad |H_3|=27,
 \label{p12-83-c20-s6:eq:heisenberg-3-nuevo}
\end{equation}
con la misma ley. El núcleo de
\(\mathcal H_9\to H_3\) es
\((3\mathbb Z/9\mathbb Z)^3\), de cardinal 27.

\subsection{Representación de Weyl--Heisenberg y fase central}

La extensión nonádica admite una realización unitaria explícita. Sea
\[
 \zeta_9:=\exp(2\pi i/9),
 \qquad
 \mathscr H_9^{\rm Sch}:=\mathbb C^9,
\]
con base \(\{|n\rangle:n\in\mathbb Z/9\mathbb Z\}\). Definimos
\begin{equation}
 X|n\rangle:=|n+1\rangle,
 \qquad
 Z|n\rangle:=\zeta_9^n|n\rangle.
 \label{p12-83-c20-s6:eq:weyl-desplazamiento-fase-u029}
\end{equation}
Se tiene
\[
 ZX=\zeta_9XZ
\]
y, por multiplicación,
\begin{equation}
 (X^aZ^b)(X^cZ^d)
 =\zeta_9^{bc}X^{a+c}Z^{b+d}.
 \label{p12-83-c20-s6:eq:cociclo-weyl-nonadico-u029}
\end{equation}

\begin{theorem}[Representación fiel de Weyl--Heisenberg]
\label{p12-83-c20-s6:thm:representacion-weyl-heisenberg-u029}
La aplicación
\begin{equation}
 \rho_{\rm WH}:\mathcal H_9\longrightarrow
 U(\mathscr H_9^{\rm Sch}),
 \qquad
 \rho_{\rm WH}(a,b,z):=\zeta_9^zX^aZ^b,
 \label{p12-83-c20-s6:eq:representacion-weyl-heisenberg-u029}
\end{equation}
es una representación unitaria fiel. Para
\(u=(a,b)\) y \(v=(c,d)\), su conmutador satisface
\begin{equation}
 [\rho_{\rm WH}(u,0),\rho_{\rm WH}(v,0)]
 =\zeta_9^{\,bc-ad}I
 =\zeta_9^{-\sigma(u,v)}I.
 \label{p12-83-c20-s6:eq:conmutador-weyl-holonomia-u029}
\end{equation}
\end{theorem}

\begin{proof}
La identidad \eqref{p12-83-c20-s6:eq:cociclo-weyl-nonadico-u029} reproduce el término
\(bc\) de \eqref{p12-83-c20-s6:eq:heisenberg-9-nuevo}; de ahí la propiedad de
homomorfismo. Si \(\zeta_9^zX^aZ^b=I\), la acción sobre la base obliga
sucesivamente a \(a=0\), \(b=0\) y \(z=0\). Esto prueba fidelidad. Comparar
los productos en los dos órdenes da
\(\zeta_9^{bc-ad}\), que es la fase alternante de
\eqref{p12-83-c20-s6:eq:forma-alternante-e6-nueva}.
\end{proof}

La representación distingue dos observables conjugados. Sean
\[
 Q_{\rm H}:=\operatorname{diag}(0,1,\ldots,8),
 \qquad
 F_{jk}:=\frac13\zeta_9^{jk},
 \qquad
 P_{\rm H}:=F^\dagger Q_{\rm H}F.
\]
Las bases propias de \(Q_{\rm H}\) y \(P_{\rm H}\) son mutuamente
insesgadas, pues \(|F_{jk}|^2=1/9\). Para todo vector unitario
\(\psi\in\mathscr H_9^{\rm Sch}\),
\begin{align}
 \operatorname{Var}_\psi(Q_{\rm H})
 \operatorname{Var}_\psi(P_{\rm H})
 &\ge
 \frac14\left|
 \langle\psi,[Q_{\rm H},P_{\rm H}]\psi\rangle
 \right|^2,
 \label{p12-83-c20-s6:eq:incertidumbre-weyl-u029}\\
 H_{Q_{\rm H}}(\psi)+H_{P_{\rm H}}(\psi)
 &\ge\log9.
 \label{p12-83-c20-s6:eq:incertidumbre-entropica-weyl-u029}
\end{align}
Estos operadores son adimensionales. Una carta física de unidades es una
realización posterior y no interviene en la extensión central.

\subsection{Acción de \(10\) desplazamientos sobre el grafo de Cayley}

Para un funcional lineal
\(\ell(a,b)=pa+qb\), ponemos
\[
 D_\ell(a,b)=\left(a,b,\frac{ab}{2}+\ell(a,b)\right),
 \qquad Z=(0,0,1),
\]
donde \(2^{-1}=2\) en \(\mathbb F_3\). El conjunto inverso-cerrado
\begin{equation}
 S_\ell=
 \{D_\ell(v):v\in\mathbb F_3^2\setminus\{0\}\}
 \cup\{Z,Z^{-1}\}
 \label{p12-83-c20-s6:eq:conjunto-conexion-e6}
\end{equation}
tiene diez elementos y no contiene la identidad.

\begin{theorem}[Independencia del calibre lineal]
\label{p12-83-c20-s6:thm:independencia-calibre-e6}
Los nueve funcionales \(\ell\) producen grafos de Cayley isomorfos. Los
nueve conjuntos \(S_\ell\) forman una sola órbita bajo
\(\operatorname{Aut}(H_3)\).
\end{theorem}

\begin{proof}
La aplicación
\[
 \phi_\ell(a,b,z)=(a,b,z+\ell(a,b))
\]
es un automorfismo central y satisface
\(\phi_\ell(S_0)=S_\ell\). Para escribir la familia completa en las
coordenadas polarizadas de \cref{p12-83-c20-s6:eq:heisenberg-9-nuevo}, sea
\(c(v,w)=b c\) para \(v=(a,b)\), \(w=(c,d)\). Dado
\(M\in\operatorname{GL}(2,3)\), el defecto
\[
 B_M(v,w)=c(Mv,Mw)-\det(M)c(v,w)
\]
es simétrico, porque las partes alternantes de los dos términos coinciden.
Definimos la corrección cuadrática
\begin{equation}
 q_M(v)=2B_M(v,v),
 \qquad
 q_M(v+w)-q_M(v)-q_M(w)=B_M(v,w),
 \label{p12-83-c20-s6:eq:correccion-cuadratica-automorfismos-h3}
\end{equation}
donde \(2=2^{-1}\) en \(\mathbb F_3\). La familia completa de
automorfismos se escribe entonces
\begin{equation}
 (v,t)\longmapsto
 \bigl(Mv,\det(M)t+q_M(v)+\ell(v)\bigr),
 \qquad M\in\operatorname{GL}(2,3),\
 \ell\in(\mathbb F_3^2)^*.
 \label{p12-83-c20-s6:eq:automorfismos-h3-completos}
\end{equation}
La identidad polar de \cref{p12-83-c20-s6:eq:correccion-cuadratica-automorfismos-h3}
prueba directamente que el producto central se conserva. La familia tiene
\(48\cdot9=432\) elementos. La enumeración exhaustiva de los
\(\binom{13}{5}=1287\) subconjuntos inverso-cerrados de cardinal diez
encuentra exactamente esos nueve con los parámetros que se derivan a
continuación.
\end{proof}

Escribamos \(\underline S=\sum_{s\in S}s\) en
\(\mathbb Z[H_3]\). El conteo de productos ordenados da
\begin{equation}
 \underline S^{\,2}
 =10e+\underline S+5(\underline{H_3}-e-\underline S)
 =5\underline{H_3}+5e-4\underline S.
 \label{p12-83-c20-s6:eq:identidad-grupo-schlafli-nueva}
\end{equation}
En efecto, la identidad recibe diez factorizaciones; un elemento de
\(S\) recibe una; cada uno de los otros dieciséis recibe cinco.

\begin{theorem}[Grafo de las veintisiete rectas]
\label{p12-83-c20-s6:thm:grafo-27-rectas-nuevo}
La matriz de adyacencia \(A_\ell\) de
\(\Gamma_\ell=\operatorname{Cay}(H_3,S_\ell)\) satisface
\begin{equation}
 A_\ell^2=5I-4A_\ell+5J.
 \label{p12-83-c20-s6:eq:identidad-schlafli-nueva}
\end{equation}
Por tanto,
\[
 \Gamma_\ell=\operatorname{srg}(27,10,1,5),
 \qquad
 \operatorname{Spec}(A_\ell)=\{10^1,1^{20},(-5)^6\}.
\]
Es isomorfo al grafo de intersección de las veintisiete rectas de una
superficie cúbica lisa.
\end{theorem}

\begin{proof}
La representación regular de
\cref{p12-83-c20-s6:eq:identidad-grupo-schlafli-nueva} produce
\cref{p12-83-c20-s6:eq:identidad-schlafli-nueva}. Sus entradas dicen que todo vértice
tiene diez vecinos, que un par adyacente tiene un vecino común y un par no
adyacente tiene cinco. Sobre la recta constante, \(A_\ell\) actúa por 10;
en su complemento, \(J=0\) y
\((A_\ell-I)(A_\ell+5I)=0\). Traza cero fija las multiplicidades 20 y 6.

En el modelo de una cúbica obtenida al soplar seis puntos de
\(\mathbb P^2\), las 27 clases son
\[
 E_i,\qquad L_{ij}=H-E_i-E_j,\qquad
 Q_i=2H-\sum_{j\ne i}E_j.
\]
Con \(H^2=1\), \(E_i^2=-1\), \(H\cdot E_i=0\) y
\(E_i\cdot E_j=0\) para \(i\ne j\), su matriz de intersección tiene los
mismos parámetros. La identificación con la configuración clásica usa el
teorema de unicidad de esta realización de las 27 rectas
\cite{Manivel2005}; no se atribuye a un certificado inexistente.
\end{proof}

La afirmación de isomorfía puede hacerse completamente efectiva mediante
una biyección calculada. El etiquetado no es absoluto: otra sección lineal
o una automorfía de la superficie produce otra tabla de la misma clase.
\begin{longtable}{@{}c@{\,}c@{\,}c@{\qquad}l@{}}
\caption{Bijección explícita entre \(H_3\) y las veintisiete clases de
rectas.}
\label{p12-83-c20-s6:tab:h3-27-rectas-u029}\\
\toprule
\(a\)&\(b\)&\(z\)&clase de recta\\
\midrule
\endfirsthead
\multicolumn{4}{@{}l}{\small\itshape Cuadro \thetable\ (continuación)}\\
\toprule
\(a\)&\(b\)&\(z\)&clase de recta\\
\midrule
\endhead
0&0&0&\(E_1\)\\
0&0&1&\(L_{12}\)\\
0&0&2&\(Q_2\)\\
0&1&0&\(L_{13}\)\\
0&1&1&\(Q_1\)\\
0&1&2&\(E_3\)\\
0&2&0&\(Q_3\)\\
0&2&1&\(E_2\)\\
0&2&2&\(L_{23}\)\\
1&0&0&\(L_{14}\)\\
1&0&1&\(L_{35}\)\\
1&0&2&\(L_{26}\)\\
1&1&0&\(L_{36}\)\\
1&1&1&\(L_{24}\)\\
1&1&2&\(L_{15}\)\\
1&2&0&\(L_{25}\)\\
1&2&1&\(L_{16}\)\\
1&2&2&\(L_{34}\)\\
2&0&0&\(Q_4\)\\
2&0&1&\(L_{46}\)\\
2&0&2&\(E_6\)\\
2&1&0&\(E_5\)\\
2&1&1&\(Q_6\)\\
2&1&2&\(L_{56}\)\\
2&2&0&\(L_{45}\)\\
2&2&1&\(E_4\)\\
2&2&2&\(Q_5\)\\
\bottomrule
\end{longtable}

Como control local, los diez vecinos de \(e=(0,0,0)\) son los elementos de
\(S_0\). La tabla los transforma en las diez clases que cortan a \(E_1\):
las cinco \(L_{1j}\) y las cinco \(Q_j\), con \(2\le j\le6\). La
verificación completa compara las 729 entradas de ambas matrices de
adyacencia.

El complemento \(B_\ell=J-I-A_\ell\) satisface
\begin{equation}
 B_\ell^2=8I+2B_\ell+8J,
 \qquad
 \operatorname{srg}(27,16,10,8),
 \qquad
 \operatorname{Spec}(B_\ell)=\{16^1,4^6,(-2)^{20}\}.
 \label{p12-83-c20-s6:eq:complemento-schlafli}
\end{equation}
Éste es el grafo de Schläfli en la convención de adyacencia complementaria;
\(A_\ell\) es el grafo de intersección de parámetros
\((27,10,1,5)\).

Cada arista pertenece a un triángulo único, pues \(\lambda=1\). El número
de triángulos es
\[
 \frac{27\cdot10}{6}=45.
\]
El grupo completo de automorfismos tiene orden
\[
 27\cdot1920=51840
\]
y, por su acción sobre la configuración de rectas, se identifica con
\(W(E_6)\), no sólo por igualdad de órdenes \cite{Manivel2005}.

\subsection{El retículo de raíces que realiza \texorpdfstring{\(E_6\)}{E6}}

Para hacer explícito el objeto sobre el que actúa ese grupo, sea
\begin{equation}
 \operatorname{Pic}(X)
 =\mathbb ZH\oplus\bigoplus_{i=1}^6\mathbb ZE_i,
 \qquad
 H^2=1,\qquad E_i^2=-1,\qquad H\cdot E_i=E_i\cdot E_j=0\ (i\ne j),
 \label{p12-83-c20-s6:eq:picard-cubica-e6}
\end{equation}
el retículo de Picard de la explosión de seis puntos generales de
\(\mathbb P^2\). Su clase canónica es
\[
 K_X=-3H+E_1+\cdots+E_6.
\]
El complemento ortogonal
\begin{equation}
 L_{E_6}=K_X^\perp\subset\operatorname{Pic}(X),
 \qquad (x,y)_{E_6}:=-x\cdot y,
 \label{p12-83-c20-s6:eq:reticulo-raices-e6}
\end{equation}
es positivo definido de rango seis. Una base de raíces simples es
\begin{equation}
 \begin{aligned}
 \alpha_1&=E_1-E_2,& \alpha_2&=E_2-E_3,&
 \alpha_3&=E_3-E_4,\\
 \alpha_4&=E_4-E_5,& \alpha_5&=E_5-E_6,&
 \alpha_6&=H-E_1-E_2-E_3.
 \end{aligned}
 \label{p12-83-c20-s6:eq:raices-simples-e6}
\end{equation}
Cada \(\alpha_i\) es ortogonal a \(K_X\), tiene norma dos y sus productos
son \(-1\) exactamente para los pares adyacentes del diagrama
\(E_6\): la cadena \(1-2-3-4-5\), con el nodo \(6\) unido al nodo \(3\).
Por tanto, su matriz de Gram es la matriz de Cartan de \(E_6\).

\begin{proposition}[Las setenta y dos raíces y sus reflexiones]
\label{p12-83-c20-s6:prop:setenta-dos-raices-e6}
Las raíces de \(L_{E_6}\) son
\begin{equation}
 \begin{aligned}
 &E_i-E_j &&(i\ne j),\\
 &\pm(H-E_i-E_j-E_k) &&(1\le i<j<k\le6),\\
 &\pm(2H-E_1-\cdots-E_6),
 \end{aligned}
 \label{p12-83-c20-s6:eq:enumeracion-raices-e6}
\end{equation}
en total \(30+40+2=72\). Para cada raíz \(\alpha\),
\begin{equation}
 s_\alpha(x)=x+(x\cdot\alpha)\alpha
 \label{p12-83-c20-s6:eq:reflexion-raiz-e6}
\end{equation}
preserva la intersección y \(K_X\). Las seis reflexiones simples generan
\(W(E_6)\) y permutan fielmente las 27 clases
\(E_i,L_{ij},Q_i\) de \cref{p12-83-c20-s6:thm:grafo-27-rectas-nuevo}.
\end{proposition}

\begin{proof}
Sustituir cada clase de \cref{p12-83-c20-s6:eq:enumeracion-raices-e6} en
\cref{p12-83-c20-s6:eq:picard-cubica-e6} da \(K_X\cdot\alpha=0\) y
\(\alpha^2=-2\). La cuenta de clases es la indicada. La fórmula de reflexión
es la reflexión ortogonal para la forma \(-(\cdot)\); preserva \(K_X\) porque
\(K_X\cdot\alpha=0\). El sistema simple de
\cref{p12-83-c20-s6:eq:raices-simples-e6} genera todas las raíces y su grupo de Coxeter es
\(W(E_6)\). Su acción sobre las 27 clases de curvas excepcionales conserva la
intersección, de modo que coincide con la acción sobre el grafo ya
construido \cite{Manivel2005}.
\end{proof}

\subsection{Estratificación de las 729 palabras}

Escribimos una palabra de seis trits como
\[
 w=(\ell_1,h_1,\ell_2,h_2,\ell_3,h_3),
\]
y definimos dos elementos de \(H_3\),
\[
 L(w)=(\ell_1,\ell_2,\ell_3),\qquad
 H(w)=(h_1,h_2,h_3),
\]
usando la ley central en sus productos. Según
\(L(w)^{-1}H(w)\), las 729 palabras se dividen en
\begin{align*}
 \mathcal D_\ell&=\{w:L(w)^{-1}H(w)=e\},\\
 \mathcal I_\ell&=\{w:L(w)^{-1}H(w)\in S_\ell\},\\
 \mathcal S_\ell&=\{w:L(w)^{-1}H(w)\notin\{e\}\cup S_\ell\}.
\end{align*}

\begin{theorem}[Partición Hensel--Schläfli]
\label{p12-83-c20-s6:thm:particion-hensel-schlafli-nueva}
\begin{equation}
 729=27+270+432,
 \label{p12-83-c20-s6:eq:particion-729-e6-nueva}
\end{equation}
correspondientes a diagonal, incidencia y alabeamiento.
\end{theorem}

\begin{proof}
Para cada uno de los 27 valores de \(L(w)\), hay un único valor de
\(H(w)\) con diferencia identidad, diez con diferencia en \(S_\ell\) y
dieciséis restantes. Se multiplica \(27\) por \(1,10,16\).
\end{proof}

El cardinal \(432\) del estrato alabeado coincide numéricamente con la
multiplicidad de la región transversal de \(\pi\) en
\cref{eq:descomposiciones-cinco-regiones-pi}. No se identifica por ello ambos objetos: aquí
se cuentan palabras clasificadas por \(L(w)^{-1}H(w)\), mientras allí se
cuentan preimágenes de un registro \(U_6\). Sin un entrelazador explícito y
equivariante entre esos dominios, \(432=432\) es sólo un control cardinal.

\subsection{Necesidad de la fase central}

La enumeración de todos los subconjuntos inverso-cerrados de cardinal diez
en los cinco grupos de orden 27 da
\[
\begin{array}{c|c|c}
\text{grupo}&\text{abeliano}&
\#\operatorname{srg}(27,10,1,5)\\ \hline
C_{27}&\text{sí}&0\\
C_9\times C_3&\text{sí}&0\\
C_3^3&\text{sí}&0\\
C_9\rtimes C_3&\text{no}&9\\
H_3&\text{no}&9.
\end{array}
\]
Así, conservar 27 etiquetas pero abelianizar la composición destruye la
configuración. No bastan la cardinalidad ni el grado.

Si \(\mathcal T\) es el conjunto de 45 triángulos, el cúbico
\[
 \mathcal C_\Gamma(x_1,\ldots,x_{27})
 =\sum_{\{i,j,k\}\in\mathcal T}x_ix_jx_k
\]
recupera la adyacencia, porque cada arista está en un triángulo único. Su
estabilizador permutacional es exactamente
\[
 \operatorname{Aut}(\Gamma_\ell)\cong W(E_6).
\]

\begin{figure}[htbp]
\centering
\[
 \text{curvaturas APP }\pm1
 \longrightarrow\sigma
 \longrightarrow\mathcal H_9
 \longrightarrow H_3
 \longrightarrow\Gamma_{27}
 \longrightarrow\mathcal T_{45}
 \longrightarrow W(E_6).
\]
\caption{Cadena causal de la rama de holonomía mixta. Es paralela a
\(C_W\to W_{12}\to M_{12}\) y no un eslabón de ella.}
\label{p12-83-c20-s6:fig:cadena-holonomia-e6-u029}
\end{figure}

\begin{criterion}[Criterios de refutación de la rama \(E_6\)]
\label{p12-83-c20-s6:crit:refutacion-rama-e6-u029}
La rama queda refutada si las curvaturas no son opuestas, si el cociclo
central se borra, si la corrección cuadrática
\eqref{p12-83-c20-s6:eq:correccion-cuadratica-automorfismos-h3} no conserva el producto,
si \eqref{p12-83-c20-s6:eq:identidad-grupo-schlafli-nueva} no produce los coeficientes
\((10,1,5)\), si un grupo abeliano produce la misma configuración, si la
partición no suma 729, si la tabla
\ref{p12-83-c20-s6:tab:h3-27-rectas-u029} no entrelaza las matrices de adyacencia o si
\(W(E_6)\) se identifica sólo por su orden y no por su acción sobre las
veintisiete clases y las 72 raíces.
\end{criterion}


\section{Acoplamiento orientado APP–Fock y rigidez transversal de \(12\) fibras}
\label{p12-83-c20-s7:chap:acoplamiento-app-fock}
\index{APP--Fock!acoplamiento orientado}
\index{Fock!grado dos}
\index{orientación de hoja}

El determinante ciclotómico ordinario no distingue \(C\) de \(C^{-1}\).
La información suma--producto se conserva introduciendo explícitamente la
paridad de hoja y una recta de orientación. El acoplamiento resultante
produce el índice \(54\) desde el agregado APP. Una vez fijadas las cuatro
condiciones del acoplamiento, \(54\), \(J\) y \(196884\) no son entradas de
la evaluación. Esta independencia evaluativa no afirma unicidad universal
del ansatz ni independencia histórica de su diseño.

\subsection{Marco triangular en cada órbita}

Sea
\[
M=\mathbb Z[\mathbb Z/9\mathbb Z]
\]
y sea \(M_{\mathrm{bal}}^{C_3}\) el submódulo de los elementos invariantes
\(\nu\) cuyos seis coeficientes unitarios coinciden:
\[
\nu_u=c(\nu)
\qquad
\bigl(u\in(\mathbb Z/9\mathbb Z)^\times\bigr).
\]
Los coeficientes de \(e_0,e_3,e_6\) permanecen libres.

Para una órbita libre
\(\mathcal O=\{u,4u,7u\}\), pongamos
\[
b_v=e_v-e_{4v}\in A(\mathcal O).
\]
Si \(b_v^\flat\) es el covector definido por la forma de \(A_2\), las tres
raíces del marco triangular satisfacen
\[
\boxed{
\sum_{v\in\mathcal O}b_v\otimes b_v^\flat
=3I_{A(\mathcal O)}.}
\]
En consecuencia,
\[
\sum_{v\in\mathcal O}\nu_v\,b_v\otimes b_v^\flat
=3c(\nu)I.
\]
El factor tres procede de la resolución de la identidad por el marco de
raíces.

\subsection{Acoplamiento orientado \(\mathcal H_{\Sigma,\Pi}^{(m)}\) entre el balance \(C_3\), la paridad de hoja y el grado \(2\) del espacio de Fock}

Para \(m\in2\mathbb N_{>0}\), sea
\[
\begin{aligned}
\mathcal V_m&=\mathcal V_{\Sigma,m}\oplus\mathcal V_{\Pi,m},\\
\mathcal V_{\Sigma,m}&=(A_2\otimes\mathbb C)^{\oplus m/2},\qquad
\mathcal V_{\Pi,m}=(A_2\otimes\mathbb C)^{\oplus m/2},\\
g_m&=C^{\oplus m/2}\oplus(C^{-1})^{\oplus m/2},\\
\mathsf E_m&=I_{\mathcal V_{\Sigma,m}}
\oplus(-I_{\mathcal V_{\Pi,m}}).
\end{aligned}
\]
La especialización \(m=12\), con las etiquetas de par, hoja y órbita,
es \(\mathcal V_{12}=W_{\mathrm{APP}}\otimes\mathbb C\). Esta notación
evita confundir el espacio representacional con el sistema de Witt
\(W_{12}=S(5,6,12)\). El grado dos del espacio de
Fock y la acción inducida son
\[
\mathcal F_2(\mathcal V_m)
=\mathcal V_m[2]\oplus\operatorname{Sym}^2(\mathcal V_m[1]),
\]
\[
\Gamma_2(g_m)
=g_m|_{\mathcal V_m[2]}
\oplus\operatorname{Sym}^2(g_m|_{\mathcal V_m[1]}).
\]

Sea \(o_{\Sigma\Pi}\) la recta de orientación con signo
\(\Pi-\Sigma\), y sea \(M_{\mathrm{bal}}^{C_3,-}\) la copia orientada en
la que el intercambio de hojas actúa por \(\nu\mapsto-\nu\). Definimos
\[
\boxed{
\begin{aligned}
\mathcal H_{\Sigma,\Pi}^{(m)}&:
M_{\mathrm{bal}}^{C_3,-}\longrightarrow
\operatorname{End}(\mathcal F_2(\mathcal V_m))\otimes o_{\Sigma\Pi},\\
\mathcal H_{\Sigma,\Pi}^{(m)}(\nu)
&=3c(\nu)
\left(0_{\mathcal V_m[2]}\oplus\operatorname{Sym}^2\mathsf E_m\right)
\otimes o_{\Sigma\Pi}.
\end{aligned}}
\]
La transformación es lineal en \(c\), está soportada en
\(\operatorname{Sym}^2(\mathcal V_m[1])\), utiliza la paridad de hoja y adopta la
normalización del marco triangular.

Sea \(\sigma\) el intercambio de hojas. Entonces
\[
\nu\longmapsto-\nu,\qquad
o_{\Sigma\Pi}\longmapsto-o_{\Sigma\Pi},\qquad
\mathsf E_m\longmapsto-\mathsf E_m,
\]
y
\[
\operatorname{Sym}^2(-\mathsf E_m)
=\operatorname{Sym}^2(\mathsf E_m).
\]
Por tanto,
\[
\mathcal H_{\Sigma,\Pi}^{(m)}(\sigma\nu)
=\sigma\mathcal H_{\Sigma,\Pi}^{(m)}(\nu).
\]
Al invertir la orientación, el covector dual cambia también de signo; así
\((-o_{\Sigma\Pi}^{\vee})(-o_{\Sigma\Pi})=1\), y la evaluación orientada
permanece normalizada.

Fijado
\(o_{\Sigma\Pi}^{\vee}(o_{\Sigma\Pi})=1\), la traza orientada es
\[
\operatorname{Tr}_{o,\mathcal F_2}
(T\otimes o_{\Sigma\Pi})
=\operatorname{Tr}_{\mathcal F_2}(T).
\]

\begin{theorem}[Índice orientado APP--Fock]
\label{p12-83-c20-s7:thm:indice-orientado-app-fock}
Supóngase que \(m\) fibras \(A_2\), con \(m\) par, se reparten por igual
entre las orientaciones \(C\) y \(C^{-1}\). Entonces
\[
\boxed{
\operatorname{Tr}_{o,\mathcal F_2}
\left(
\mathcal H_{\Sigma,\Pi}^{(m)}(\nu)\Gamma_2(g_m)
\right)
=-\frac{3m\,c(\nu)}2.}
\]
Para el agregado producido por APP,
\[
\delta=\nu_\Pi-\nu_\Sigma
=12e_0+3e_3+3e_6
-3\sum_{u\in(\mathbb Z/9\mathbb Z)^\times}e_u,
\]
se tiene \(c(\delta)=-3\) y
\[
m=3\ \text{pares}\cdot2\ \text{hojas}\cdot2\ \text{órbitas}=12.
\]
Por tanto,
\[
\boxed{
\operatorname{Tr}_{o,\mathcal F_2}
\left(
\mathcal H_{\Sigma,\Pi}^{(12)}(\delta)\Gamma_2(g_{12})
\right)=54.}
\]
\end{theorem}

\begin{proof}
Sea \(A=\mathsf E_mg_m\). Las contribuciones de las dos hojas se cancelan:
\(\operatorname{tr}A=0\). Como \(\mathsf E_m\) conmuta con \(g_m\) y
\(\mathsf E_m^2=I\),
\[
\operatorname{tr}(A^2)=\operatorname{tr}(g_m^2)=-m.
\]
La identidad
\[
\operatorname{tr}(\operatorname{Sym}^2A)
=\frac{(\operatorname{tr}A)^2+\operatorname{tr}(A^2)}2
\]
da \(-m/2\); el marco aporta \(3c(\nu)\).
\end{proof}

\begin{corollary}[Minimalidad del grado bosónico detector]
\label{p12-83-c20-s7:cor:minimalidad-grado-dos-app-fock-u030}
En la clase de acoplamientos declarada, el grado uno no detecta el agregado
orientado y el grado dos es el primer grado bosónico positivo que lo
detecta:
\[
 \operatorname{tr}(\mathsf E_mg_m)=0,
 \qquad
 \operatorname{tr}\!\left(
 \operatorname{Sym}^2(\mathsf E_mg_m)
 \right)=-\frac m2\neq0.
\]
\end{corollary}

\begin{proof}
La primera igualdad es la cancelación de las dos hojas. La segunda se
obtuvo en la prueba del \cref{p12-83-c20-s7:thm:indice-orientado-app-fock}. Como el grado
cero sólo contiene la unidad y el grado uno da traza nula, el grado dos es
el primero con respuesta no nula.
\end{proof}

\subsection{Rigidez transversal del número de fibras}

Para un entero \(m\ge1\), definimos la serie formal
\begin{equation}
 T_m(q):=
 q^{-1}\prod_{n\ge1}(1+q^n+q^{2n})^{-m}
 \in q^{-1}\mathbb Z[[q]]
 \label{p12-83-c20-s7:eq:familia-fock-m-fibras-u030}
\end{equation}
y, en la suma ortogonal \(A_2^m\), el vector radial
\begin{equation}
 v_m:=4\rho_1+\rho_2+\cdots+\rho_m,
 \qquad
 \|\rho_j\|^2=2.
 \label{p12-83-c20-s7:eq:familia-radial-m-fibras-u030}
\end{equation}
Construimos dos funciones enteras sin usar previamente \(m=12\):
\[
 F(m):=[q]T_m(q),
 \qquad
 R(m):=\|v_m\|^2.
\]

\begin{proposition}[Coeficiente de Fock y norma radial]
\label{p12-83-c20-s7:prop:dos-funciones-rigidez-doce-u030}
Para todo \(m\ge1\),
\begin{equation}
 F(m)=\frac{m(m-3)}2,
 \qquad
 R(m)=2(m+15).
 \label{p12-83-c20-s7:eq:funciones-rigidez-doce-u030}
\end{equation}
\end{proposition}

\begin{proof}
Hasta orden dos,
\[
 (1+q+q^2)^{-m}
 =1-mq+\frac{m(m-1)}2q^2+O(q^3),
\]
\[
 (1+q^2+q^4)^{-m}=1-mq^2+O(q^4).
\]
Los factores con \(n\ge3\) no contribuyen al coeficiente de \(q^2\).
Como el factor exterior de \(T_m\) es \(q^{-1}\),
\[
 F(m)=\frac{m(m-1)}2-m=\frac{m(m-3)}2.
\]
Por ortogonalidad de las \(m\) copias de \(A_2\),
\[
 R(m)=4^2\|\rho_1\|^2+
 \sum_{j=2}^{m}\|\rho_j\|^2
 =32+2(m-1)=2(m+15).
\]
\end{proof}

\begin{theorem}[Rigidez ciclotómica--radial]
\label{p12-83-c20-s7:thm:rigidez-transversal-doce-u030}
La igualdad \(F(m)=R(m)\) equivale a
\begin{equation}
 (m-12)(m+5)=0.
 \label{p12-83-c20-s7:eq:ecuacion-rigidez-doce-u030}
\end{equation}
En el dominio \(m\in\mathbb Z_{>0}\), la única solución es
\[
 \boxed{m=12,\qquad F(12)=R(12)=54.}
\]
\end{theorem}

\begin{proof}
Igualar las dos expresiones de
\eqref{p12-83-c20-s7:eq:funciones-rigidez-doce-u030} produce
\[
 m(m-3)=4(m+15),
\]
es decir,
\[
 m^2-7m-60=(m-12)(m+5)=0.
\]
La segunda raíz es negativa y queda fuera del dominio.
\end{proof}

Esta rigidez pertenece a la familia ciclotómica--radial declarada. No usa el
valor \(54\) para escoger \(m\): construye primero \(F(m)\) y \(R(m)\), y
resuelve después su igualdad.

\subsection{Factorización del acoplamiento APP–Fock: hipótesis, invariantes y condiciones necesarias de validez y obstrucciones}

El acoplamiento factoriza por
\[
M_{\mathrm{bal}}^{C_3,-}/\ker c
\cong\mathbb Z\otimes o_{\Sigma\Pi},
\]
y el mapa inducido es inyectivo porque
\(\operatorname{Sym}^2\mathsf E_m\neq0\).
La fórmula queda seleccionada por las cuatro condiciones declaradas:
linealidad en \(c\), soporte en
\(\operatorname{Sym}^2(\mathcal V_m[1])\), paridad
de hoja y normalización por el marco triangular. No se afirma unicidad
entre todos los acoplamientos naturales ni independencia histórica del
ansatz. Una vez fijadas esas condiciones, \(54\), \(J\) y \(196884\) no
son entradas de la evaluación.

Tres controles muestran que el resultado no está fijado por una
normalización retrospectiva:
\[
(m,c)=(12,-2)\longmapsto36,\qquad
(m,c)=(10,-3)\longmapsto45,
\]
y
\[
\delta_k=\delta+ke_0
\]
deja invariante \(\mathcal H_{\Sigma,\Pi}^{(12)}\) mientras el peso digital
cambia de \(54\) a \(54+9k\). En particular, para \(k=1\),
\[
\operatorname{Tr}_{o,\mathcal F_2}
\left(
\mathcal H_{\Sigma,\Pi}^{(12)}(\delta+e_0)\Gamma_2(g_{12})
\right)=54,
\qquad
w(\delta+e_0)=63.
\]

\begin{proposition}[Necesidad del dato de orientación]
\label{p12-83-c20-s7:prop:orientacion-app-fock}
Si se olvida \(o_{\Sigma\Pi}\), todo escalar natural lineal en
\(\nu_\Pi-\nu_\Sigma\) y dependiente sólo de la clase de la representación
se anula.
\end{proposition}

\begin{proof}
\(C\) y \(C^{-1}\) son conjugados, por lo que la clase representacional es
par bajo el intercambio de hojas. El agregado
\(\nu_\Pi-\nu_\Sigma\) es impar. Sin la recta de orientación, una
transformación simultáneamente par e impar vale cero.
\end{proof}

El teorema establece así el transporte exacto del agregado al grado dos de
Fock. El defecto local estado a estado conserva la obstrucción de descenso
del \cref{p12-83-c1-s5:cor:obstruccion-isotipica-app-autonoma}; ambos enunciados pertenecen a
dominios distintos y son compatibles.


% Unidad U031. Alineamiento APP--Witt, automorfia de orden tres,
% construccion VOA y Moonshine.
